% =====================================================================
%  Real-Time Hand & Finger Tracker with MediaPipe and OpenCV
%  A Complete Study Handbook
%
%  Build (from this directory):
%      pdflatex Study-Handbook.tex
%      pdflatex Study-Handbook.tex      <- twice, for the table of contents
%
%  The chapter bodies live in ch01..ch05 + appendix-source.tex.
%  Chapter 5 and the appendices quote the real project source files with
%  \lstinputlisting, so the code printed here cannot drift from the code
%  that ships. Build from the doc/ directory.
% =====================================================================

\documentclass[11pt,a4paper,oneside]{report}

% ---------- Encoding, fonts, typography ------------------------------
\usepackage[T1]{fontenc}
\usepackage[utf8]{inputenc}
\usepackage{lmodern}
\usepackage[scaled=0.90]{helvet}
\usepackage{inconsolata}
\usepackage{textcomp}
\usepackage{microtype}
\usepackage[left=1in,right=1in,top=1.05in,bottom=1in]{geometry}
\usepackage{amsmath,amssymb}
\usepackage{graphicx}
\usepackage[table]{xcolor}
\usepackage{booktabs,array,tabularx,longtable}
\usepackage{enumitem}
\usepackage{caption}
\usepackage{needspace}
\usepackage{emptypage}

% ---------- Palette --------------------------------------------------
\definecolor{navy}{HTML}{1E293B}        % primary headers / titles
\definecolor{navydark}{HTML}{0F172A}    % code panel background
\definecolor{cyanx}{HTML}{0EA5E9}       % primary accent
\definecolor{tealx}{HTML}{0D9488}       % secondary accent
\definecolor{amberx}{HTML}{F59E0B}      % warnings
\definecolor{redx}{HTML}{EF4444}        % traps / errors
\definecolor{violetx}{HTML}{8B5CF6}     % ML cards
\definecolor{slatebg}{HTML}{F8FAFC}     % box backgrounds
\definecolor{slateline}{HTML}{E2E8F0}   % table striping / rules
\definecolor{slatetext}{HTML}{475569}   % muted body text
\definecolor{skysoft}{HTML}{EFF6FF}     % python card background
\definecolor{violetsoft}{HTML}{F5F3FF}  % ml card background
\definecolor{tealsoft}{HTML}{F0FDFA}    % math card background
\definecolor{ambersoft}{HTML}{FFFBEB}   % trap card background
\definecolor{codekey}{HTML}{7DD3FC}     % listings: keywords
\definecolor{codestr}{HTML}{FCD34D}     % listings: strings
\definecolor{codecmt}{HTML}{94A3B8}     % listings: comments
\definecolor{codenum}{HTML}{FCA5A5}     % listings: numbers
\definecolor{codetext}{HTML}{E2E8F0}    % listings: body text

% ---------- Headers and footers --------------------------------------
\usepackage{fancyhdr}
\pagestyle{fancy}
\fancyhf{}
\renewcommand{\headrulewidth}{0.4pt}
\renewcommand{\footrulewidth}{0pt}
\renewcommand{\headrule}{\hbox to\headwidth{\color{slateline}\leaders\hrule height \headrulewidth\hfill}}
\fancyhead[L]{\small\sffamily\color{slatetext}\nouppercase{\leftmark}}
\fancyhead[R]{\small\sffamily\color{navy}\thepage}
\fancypagestyle{plain}{%
  \fancyhf{}%
  \renewcommand{\headrulewidth}{0pt}%
  \fancyhead[R]{\small\sffamily\color{navy}\thepage}%
}

% ---------- Chapter and section styling ------------------------------
\usepackage{titlesec}
\titleformat{\chapter}[display]
  {\normalfont\sffamily\bfseries\color{navy}}
  {\flushright\large\sffamily\color{cyanx}\MakeUppercase{\chaptertitlename}\ \thechapter}
  {0.6em}
  {\Huge\flushright}
  [\vspace{0.4em}{\color{slateline}\titlerule[1.2pt]}]
\titlespacing*{\chapter}{0pt}{-18pt}{22pt}

\titleformat{\section}
  {\normalfont\Large\sffamily\bfseries\color{navy}}{\color{tealx}\thesection}{0.6em}{}
\titleformat{\subsection}
  {\normalfont\large\sffamily\bfseries\color{navy}}{\color{cyanx}\thesubsection}{0.6em}{}
\titleformat{\subsubsection}
  {\normalfont\normalsize\sffamily\bfseries\color{slatetext}}{\thesubsubsection}{0.6em}{}
\titlespacing*{\section}{0pt}{1.6em}{0.5em}
\titlespacing*{\subsection}{0pt}{1.2em}{0.4em}

% ---------- Table of contents ----------------------------------------
\usepackage{tocloft}
\renewcommand{\cfttoctitlefont}{\Huge\sffamily\bfseries\color{navy}}
\renewcommand{\cftaftertoctitle}{\vspace{0.4em}{\color{slateline}\hrule height 1.2pt}\vspace{0.6em}}
\renewcommand{\cftchapfont}{\sffamily\bfseries\color{navy}}
\renewcommand{\cftchappagefont}{\sffamily\bfseries\color{cyanx}}
\renewcommand{\cftsecfont}{\color{navy}}
\renewcommand{\cftsecpagefont}{\color{tealx}}
\renewcommand{\cftsubsecfont}{\small\color{slatetext}}
\renewcommand{\cftsubsecpagefont}{\small\color{slatetext}}
\renewcommand{\cftdotsep}{1.2}
\setlength{\cftbeforechapskip}{0.7em}
\setlength{\cftbeforesecskip}{0.25em}

% ---------- Links -----------------------------------------------------
\usepackage[colorlinks=true,linkcolor=cyanx,citecolor=tealx,urlcolor=cyanx,
            pdftitle={Real-Time Hand and Finger Tracker: A Complete Study Handbook},
            pdfauthor={Project Study Handbook}]{hyperref}

% ---------- Diagrams and teaching boxes ------------------------------
\usepackage[many,skins,breakable]{tcolorbox}
\usepackage{tikz}
\usetikzlibrary{positioning,arrows.meta,shapes.geometric,calc,fit,backgrounds,
                decorations.pathreplacing,decorations.markings}

% The 21 MediaPipe hand landmarks, laid out for drawing: a right hand,
% presented palm-to-camera, fingers up, in TikZ coordinates (y grows upward).
% Defined once here and reused by the title page and the Chapter 3 diagrams
% so every hand drawing in this handbook has identical geometry.
\newcommand{\handcoordinates}{%
  \coordinate (L0)  at ( 0.00, 0.00);   % wrist
  \coordinate (L1)  at ( 1.10,-0.35);   % thumb CMC
  \coordinate (L2)  at ( 1.90, 0.15);   % thumb MCP
  \coordinate (L3)  at ( 2.45, 0.85);   % thumb IP
  \coordinate (L4)  at ( 2.75, 1.55);   % thumb tip
  \coordinate (L5)  at ( 0.85, 1.15);   % index MCP
  \coordinate (L6)  at ( 0.95, 2.10);   % index PIP
  \coordinate (L7)  at ( 1.00, 2.60);   % index DIP
  \coordinate (L8)  at ( 1.05, 3.05);   % index tip
  \coordinate (L9)  at ( 0.15, 1.35);   % middle MCP
  \coordinate (L10) at ( 0.15, 2.35);   % middle PIP
  \coordinate (L11) at ( 0.15, 2.90);   % middle DIP
  \coordinate (L12) at ( 0.15, 3.40);   % middle tip
  \coordinate (L13) at (-0.55, 1.25);   % ring MCP
  \coordinate (L14) at (-0.60, 2.20);   % ring PIP
  \coordinate (L15) at (-0.62, 2.70);   % ring DIP
  \coordinate (L16) at (-0.65, 3.15);   % ring tip
  \coordinate (L17) at (-1.15, 1.00);   % pinky MCP
  \coordinate (L18) at (-1.30, 1.75);   % pinky PIP
  \coordinate (L19) at (-1.38, 2.15);   % pinky DIP
  \coordinate (L20) at (-1.45, 2.50);   % pinky tip
}
% The 21 skeletal connections MediaPipe reports as HAND_CONNECTIONS.
\newcommand{\handbones}{%
  \foreach \a/\b in {0/1,1/2,2/3,3/4,0/5,5/6,6/7,7/8,5/9,9/10,10/11,11/12,
                     9/13,13/14,14/15,15/16,13/17,17/18,18/19,19/20,0/17}{%
    \draw[handbone] (L\a) -- (L\b);}%
}
\tikzset{
  handbone/.style={draw=slateline, line width=1.1pt},
  joint/.style={circle, fill=cyanx, inner sep=0pt, minimum size=4.4pt},
  jointtip/.style={circle, fill=amberx, inner sep=0pt, minimum size=6pt},
  jointindex/.style={circle, fill=redx, inner sep=0pt, minimum size=6.8pt},
}

% ---------- Code listings --------------------------------------------
\usepackage{listings}
\lstdefinestyle{py}{
  language=Python,
  basicstyle=\ttfamily\footnotesize\color{codetext},
  keywordstyle=\color{codekey}\bfseries,
  commentstyle=\color{codecmt}\itshape,
  stringstyle=\color{codestr},
  numberstyle=\ttfamily\tiny\color{codecmt},
  numbers=left, numbersep=9pt,
  backgroundcolor=\color{navydark},
  showstringspaces=false,
  breaklines=true, breakatwhitespace=false,
  tabsize=4, upquote=true,
  columns=fullflexible, keepspaces=true,
  xleftmargin=1.4em, xrightmargin=0.4em, framexleftmargin=1.2em,
  extendedchars=true,
  literate={---}{{\textcolor{codecmt}{---}}}1
           {--}{{\textcolor{codecmt}{--}}}1
}
\lstdefinestyle{codetextonly}{
  language=,
  basicstyle=\ttfamily\footnotesize\color{codetext},
  backgroundcolor=\color{navydark},
  numbers=none, showstringspaces=false, breaklines=true,
  columns=fullflexible, keepspaces=true,
  xleftmargin=0.8em, xrightmargin=0.4em,
}
\lstset{style=py}

% Rounded dark card that wraps a listing. Breakable, because several quoted
% blocks are longer than a page's remaining space, and an unbreakable box would
% push them to the next page and leave a large gap behind.
\newtcolorbox{codecard}[1][]{
  enhanced, breakable=true,
  colback=navydark, colframe=navydark,
  boxrule=0pt, arc=1.5mm,
  left=0.5mm, right=0.5mm, top=1mm, bottom=1mm,
  #1
}

% Inline code: \code{a_b} prints a_b literally, no escaping needed.
\newcommand{\code}[1]{\texttt{\small\detokenize{#1}}}

% ---------- Teaching environments ------------------------------------
\newtcolorbox[auto counter, number within=chapter]{pythonconcept}[2][]{%
  enhanced, breakable,
  colback=skysoft, colframe=cyanx, boxrule=0.8pt, arc=2mm,
  left=3mm, right=3mm, top=2.5mm, bottom=2.5mm,
  fonttitle=\sffamily\bfseries,
  coltitle=white, colbacktitle=cyanx,
  title={Python Concept \thetcbcounter\ \textbullet\ #2},
  attach boxed title to top left={xshift=3mm, yshift=-2.5mm},
  boxed title style={arc=1mm, boxrule=0pt},
  #1}

\newtcolorbox[auto counter, number within=chapter]{mlconcept}[2][]{%
  enhanced, breakable,
  colback=violetsoft, colframe=violetx, boxrule=0.8pt, arc=2mm,
  left=3mm, right=3mm, top=2.5mm, bottom=2.5mm,
  fonttitle=\sffamily\bfseries,
  coltitle=white, colbacktitle=violetx,
  title={ML Concept \thetcbcounter\ \textbullet\ #2},
  attach boxed title to top left={xshift=3mm, yshift=-2.5mm},
  boxed title style={arc=1mm, boxrule=0pt},
  #1}

\newtcolorbox[auto counter, number within=chapter]{mathdeepdive}[2][]{%
  enhanced, breakable,
  colback=tealsoft, colframe=tealx, boxrule=1pt, arc=1mm,
  left=3mm, right=3mm, top=2.5mm, bottom=2.5mm,
  fonttitle=\sffamily\bfseries,
  coltitle=white, colbacktitle=tealx,
  title={Math Deep Dive \thetcbcounter\ \textbullet\ #2},
  attach boxed title to top left={xshift=3mm, yshift=-2.5mm},
  boxed title style={arc=1mm, boxrule=0pt},
  #1}

\newtcolorbox[auto counter, number within=chapter]{edgetrap}[2][]{%
  enhanced, breakable,
  colback=ambersoft, colframe=amberx, boxrule=1pt, arc=2mm,
  left=3mm, right=3mm, top=2.5mm, bottom=2.5mm,
  fonttitle=\sffamily\bfseries,
  coltitle=white, colbacktitle=amberx,
  title={Edge-Case Trap \thetcbcounter\ \textbullet\ #2},
  attach boxed title to top left={xshift=3mm, yshift=-2.5mm},
  boxed title style={arc=1mm, boxrule=0pt},
  borderline west={2.5pt}{0pt}{redx},
  #1}

% Small inline "key idea" strip.
\newtcolorbox{keyidea}[1][]{%
  enhanced, breakable,
  colback=slatebg, colframe=slateline, boxrule=0.6pt, arc=1mm,
  left=3mm, right=3mm, top=2mm, bottom=2mm, #1}

\newcommand{\term}[1]{\textbf{\color{navy}#1}}

% ---------- Lists -----------------------------------------------------
\setlist[itemize]{leftmargin=1.4em, itemsep=0.25em, topsep=0.4em}
\setlist[enumerate]{leftmargin=1.6em, itemsep=0.3em, topsep=0.4em}
\setlist[description]{leftmargin=1.4em, itemsep=0.25em}

% ---------- Captions --------------------------------------------------
\captionsetup{font={small,sf},labelfont={bf,color=navy},labelsep=period}

% Longer identifiers (\code{...}, file paths) cannot hyphenate, so give TeX
% extra room to stretch spaces rather than run into the margin.
\emergencystretch=2.5em
\setlength{\parskip}{0pt}

% =====================================================================
\begin{document}

% ---------- Title page -----------------------------------------------
\thispagestyle{empty}
\begin{tikzpicture}[remember picture, overlay]
  \fill[navy] (current page.north west) rectangle ([yshift=-9.9cm]current page.north east);
  \fill[cyanx] ([yshift=-9.9cm]current page.north west) rectangle ([yshift=-10.15cm]current page.north east);
  \fill[tealx] ([yshift=-10.15cm]current page.north west) rectangle ([yshift=-10.25cm]current page.north east);
\end{tikzpicture}

\vspace*{1.1cm}
\begin{flushleft}
  {\sffamily\bfseries\color{white}\fontsize{30}{34}\selectfont Real-Time Hand \&\\[2pt]
   Finger Tracker\par}
  \vspace{0.7cm}
  {\sffamily\color{cyanx}\Large with MediaPipe and OpenCV\par}
  \vspace{0.35cm}
  {\sffamily\color{slateline}\large A Complete Study Handbook\par}
\end{flushleft}

\vspace{4.6cm}

\begin{center}
\begin{tikzpicture}[scale=0.72, every node/.style={font=\scriptsize\sffamily}]
  \handcoordinates
  \handbones
  \foreach \i in {0,...,20} {\node[joint] at (L\i) {};}
  \foreach \i in {4,12,16,20} {\node[jointtip] at (L\i) {};}
  \node[jointindex] at (L8) {};
  \node[below, color=slatetext] at (L0) {0 wrist};
  \node[right, color=slatetext] at (L4) {4 thumb tip};
  \node[above right, color=slatetext] at (L8) {8 index tip};
  \node[left, color=slatetext] at (L17) {17 pinky MCP};
\end{tikzpicture}
\end{center}

\vfill
\begin{flushleft}
  {\sffamily\color{slatetext}\small
  Python 3.9--3.12 \textbullet\ MediaPipe 0.10.14 \textbullet\ OpenCV 4.11.0.86 \textbullet\ NumPy $<2$\\[4pt]
  Everything in this handbook is traced to real, runnable code:\\
  \code{hand_logic.py}, \code{hand_tracker.py}, \code{tests/test_hand_logic.py}\par}
\end{flushleft}
\vspace{0.6cm}

\clearpage

% ---------- How to use / legend --------------------------------------
\chapter*{How to Use This Handbook}
\addcontentsline{toc}{chapter}{How to Use This Handbook}
\markboth{How to Use This Handbook}{}

This handbook assumes \textbf{almost no Python} and \textbf{no machine learning at all}. Every idea
is introduced with a physical analogy first, then made exact with mathematics, then connected to the
actual line of code that implements it. Nothing is hand-waved and nothing is left as an exercise.

\section*{Reading order}

\begin{enumerate}
  \item \textbf{Chapters 1--2} build the two backgrounds you are missing: the Python language, how
        images are stored in memory, what a neural network is, and how MediaPipe actually finds a
        hand. No geometry yet.
  \item \textbf{Chapter 3} is the mathematical heart: the inverted $y$-axis, the thumb rules, the
        palm normal, and the two signal-processing filters.
  \item \textbf{Chapter 4} is system design: why the project is split in two modules, how the
        per-frame cache works, and how the app survives a missing camera.
  \item \textbf{Chapter 5} walks the source code block by block, quoting the real files.
  \item \textbf{Appendices} hold the complete verbatim source, the verification checklist, and a
        glossary.
\end{enumerate}

You do \emph{not} need a webcam to read any of this, and you do not need a webcam to run most of the
code: the project ships a headless \code{--image} mode that runs the entire pipeline on a still
picture.

\section*{The four teaching cards}

Throughout the handbook you will meet four coloured boxes. They always mean the same thing:

\begin{pythonconcept}{Python language, from zero}
  Soft blue cards teach the Python language itself: what a variable is, what \code{self} means, why
  a \code{deque} is not a list. If you already know Python, you can skim these --- but they also
  explain the \emph{idiom}, not just the syntax.
\end{pythonconcept}

\begin{mlconcept}{Computer vision and machine learning}
  Soft purple cards demystify the ML half: what a tensor is, what a convolutional layer does, what
  ``inference'' means. No mathematics of training is assumed or required.
\end{mlconcept}

\begin{mathdeepdive}{Derivations step by step}
  Teal framed cards derive the geometry and signal processing used by the app, with every algebraic
  step shown and a worked numeric example. Nothing is asserted without proof.
\end{mathdeepdive}

\begin{edgetrap}{Physical failure modes}
  Amber cards with a red spine flag the ways a hand tracker breaks in the real world: palm
  inversion, the adducted ``knife hand'', camera-package conflicts. These are the difference
  between a demo and a tool.
\end{edgetrap}

\section*{How code is presented}

Code in this handbook is \textbf{quoted directly from the project files}, not retyped. That means
the line numbers printed in the margin are the real line numbers of the real file, so you can open
\code{hand_logic.py}, jump to line 127, and be looking at exactly what is printed here. Chapter 5
quotes ranges; the appendices print each file whole.

\begin{keyidea}
\textbf{A note on the printed colours.} The listing panels use the project's own palette --- midnight
navy background, cyan keywords, amber strings. If you are reading a greyscale printout, keywords
appear bold, strings appear in a lighter grey, and comments are italic.
\end{keyidea}

\clearpage

% ---------- Table of contents ----------------------------------------
\renewcommand{\contentsname}{Contents}
\tableofcontents

\clearpage

% ---------- Chapters --------------------------------------------------
\chapter{Foundations for Absolute Beginners}
\label{ch:foundations}

This chapter builds the two backgrounds you need before any hand tracking can make sense: the
Python language (Section~\ref{sec:python}) and how pictures live inside a computer
(Section~\ref{sec:images}). It closes with a plain-English account of what machine learning actually
is (Section~\ref{sec:ml}). Nothing here is specific to hands yet --- that is deliberate. If you have
never written a line of code, read every word. If you have written Python before, read the boxes
anyway: they explain the \emph{why} behind idioms you may already use by habit.

\section{Python from Scratch for Computer Vision}
\label{sec:python}

\subsection{What is an interpreter, and how does Python run a file?}

When you type \code{python hand_tracker.py} and press Enter, you are not ``starting a program'' in
the way you start a compiled \code{.exe}. You are starting a program called the \emph{interpreter},
and handing it a text file.

\begin{pythonconcept}{The interpreter, and why it matters here}
An \textbf{interpreter} is a program that reads your text file and performs the instructions one
statement at a time, top to bottom, translating as it goes.

A useful analogy: a \textbf{compiled} language is like translating an entire novel into another
language and then publishing it --- slow to prepare, fast to read. An \textbf{interpreted} language
is like a simultaneous translator standing next to you --- ready instantly, but working sentence by
sentence.

Two practical consequences you will meet within the first five minutes:

\begin{itemize}
  \item \textbf{Errors appear only when that line is reached.} If line 300 has a typo but your
        camera never opens, the program may die at line 40 and you will never see the line-300
        problem. An interpreter does not read the whole file first.
  \item \textbf{The same file can behave differently on different runs} if it depends on something
        outside the file (a camera, a network, a random number). This is why the tracker logs
        ``Camera 0 opened at 640x480'' --- so you know what the run actually did.
\end{itemize}

Python files are named with the \code{.py} suffix, and \code{python} is the interpreter. In this
project there are two of them: \code{hand_logic.py} and \code{hand_tracker.py}.
\end{pythonconcept}

\subsection{Variables and dynamic typing}

A variable is a name attached to a value. That is the whole idea.

\begin{codecard}
\begin{lstlisting}
fps = 28            # an int  (whole number)
alpha = 0.65        # a float (decimal number)
label = "Right"     # a str  (text)
is_up = True        # a bool (True / False)
\end{lstlisting}
\end{codecard}

These four lines use the four basic types you will meet constantly in this project. Python calls
its basic types \emph{built-in types}, and the project uses exactly these:

\begin{center}
\rowcolors{2}{slatebg}{white}
\begin{tabularx}{\textwidth}{@{}l l X@{}}
\toprule
\rowcolor{white}
\textbf{Type} & \textbf{Example} & \textbf{Used in the project for} \\
\midrule
\code{int}   & \code{320}      & pixel coordinates, frame counters, finger counts \\
\code{float} & \code{0.65}     & normalized coordinates, the palm normal $Z$, alpha values \\
\code{str}   & \code{"Right"}  & handedness labels, gesture names, file paths \\
\code{bool}  & \code{True}     & ``is this finger extended?'' --- five per hand, every frame \\
\bottomrule
\end{tabularx}
\end{center}

\begin{pythonconcept}{Dynamic typing: the name has a type, not the box}
In C or Java you would write \code{int fps = 28;} --- the \emph{container} is declared as an integer
and can never hold anything else.

Python is \textbf{dynamically typed}: the \emph{name} simply points at a value, and the value knows
its own type. This line is legal:

\begin{lstlisting}
value = 5          # value points at an int
value = "hello"    # now the same name points at a string
\end{lstlisting}

The price of this freedom is that type mistakes are discovered at \emph{run time}, not
compile time. The project pays that price back with \textbf{type hints} and \textbf{tests} --- see
Section~\ref{sec:typehints}.
\end{pythonconcept}

\subsection{Lists, tuples and dictionaries}

Nearly every bug in a hand tracker is really a data-shape bug: the wrong number of points, points in
the wrong order, or a name that does not exist. So it is worth being precise about the three
containers this project uses.

\subsubsection*{Lists --- ordered, changeable}

A \textbf{list} is an ordered sequence you can modify. Square brackets, comma separated.

\begin{codecard}
\begin{lstlisting}
FINGER_PAIRS = [(8, 6), (12, 10), (16, 14), (20, 18)]
states = [False] * 5          # five False values
states[0] = True              # now [True, False, False, False, False]
print(len(states))            # 5
\end{lstlisting}
\end{codecard}

Line 1 is taken from \code{hand_logic.py}: it is the list of (fingertip, knuckle) landmark index
pairs for the four long fingers. Line 2 is the ``safe result'' the project falls back to whenever
landmarks are missing --- note that multiplying a list by an integer \emph{repeats} it, it does not do
arithmetic.

The most important property of a list here is \textbf{indexing}: \code{points[0]} is the first
element, \code{points[8]} is the ninth. In this project, landmark indices are the entire language we
speak --- index 0 is the wrist, index 8 is the tip of the index finger, index 17 is the base of the
pinky. If you remember one thing from this section, remember that \emph{the index number is a name}.

\subsubsection*{Tuples --- ordered, frozen}

A \textbf{tuple} is like a list that cannot be changed. Round brackets.

\begin{codecard}
\begin{lstlisting}
tip = (320, 240)          # a single point: x, y
x, y = tip                # unpacking: x becomes 320, y becomes 240
pixel = points[8]         # a tuple from the landmark list
dist = math.hypot(pixel[0] - other[0], pixel[1] - other[1])
\end{lstlisting}
\end{codecard}

A 2D point is naturally a tuple: it is a fixed pair, and it should not be accidentally mutated. The
project uses tuples for coordinates and lists for collections of coordinates, so the shape of the
data tells you what it is.

\subsubsection*{Dictionaries --- named, lookup by key}

A \textbf{dictionary} maps keys to values. Curly braces.

\begin{codecard}
\begin{lstlisting}
hand = {
    "label": "Right",
    "px": [(320, 420), (365, 400)],   # pixel landmarks
    "palm_z": -8775.0,
    "orientation": "Palm",
}

print(hand["label"])        # Right
print(hand["bbox"])         # KeyError!  this key does not exist
\end{lstlisting}
\end{codecard}

This is precisely the structure \code{parse_hand_data} builds once per frame for each detected hand.
A dictionary is the right choice because a hand has \emph{named} properties whose order does not
matter --- and because a reader of the code can ask for \code{hand["orientation"]} instead of
remembering ``the fifth element''.

\begin{edgetrap}{AttributeError and KeyError --- the two errors of this section}
\begin{itemize}
  \item \code{hand["bboxx"]} raises \code{KeyError}: the name is not in the dictionary.
  \item \code{hand.label} raises \code{AttributeError}: dictionaries are indexed with
        \code{[]}, not with a dot. The dot syntax is for \emph{objects} (Section~\ref{sec:oop}).
\end{itemize}
These two messages account for a large share of real-world Python errors in computer-vision code,
because MediaPipe results are objects while our cached hand data is a dictionary --- the access
syntax differs between them.
\end{edgetrap}

\subsection{Functions, arguments, return values and type hints}
\label{sec:typehints}

A \textbf{function} is a named block of work with inputs and (usually) one output.

\begin{codecard}
\begin{lstlisting}
def euclidean_distance(a, b):
    return math.hypot(a[0] - b[0], a[1] - b[1])

gap = euclidean_distance((0, 0), (3, 4))    # gap is 5.0
\end{lstlisting}
\end{codecard}

Reading that from left to right:

\begin{itemize}
  \item \code{def} announces a function definition.
  \item \code{euclidean\_distance} is its name.
  \item \code{(a, b)} are its \textbf{parameters} --- placeholders for the values it will receive.
  \item The \textbf{indented block} is the \textbf{body}.
  \item \code{return} hands a value back to whoever called it. A function without a \code{return}
        hands back a special value called \code{None}, meaning ``nothing''.
\end{itemize}

\subsubsection*{Default arguments}

\begin{codecard}
\begin{lstlisting}
def fingers_state(pixels, hand_label="Unknown", thumb_mode="distance"):
    ...

fingers_state(points)                       # label is "Unknown", mode "distance"
fingers_state(points, "Right")              # label is "Right"
fingers_state(points, thumb_mode="x")       # keyword argument: order does not matter
\end{lstlisting}
\end{codecard}

An argument with an \code{=} in the definition is \textbf{optional}. This is why the project's
public methods can be called simply --- \code{tracker.find\_positions(frame, 0)} --- while still
offering knobs such as \code{alpha=0.65}.

\subsubsection*{Type hints are documentation that can be checked}

\begin{codecard}
\begin{lstlisting}
def to_pixel_landmarks(landmarks, width, height):
    """Convert normalized landmarks to integer pixel coordinates."""
    return [(int(lm.x * width), int(lm.y * height)) for lm in landmarks]
\end{lstlisting}
\end{codecard}

Python does not enforce the types written in a hint --- it will happily run and crash later if you
pass nonsense. The hint is a promise to the reader, and modern editors check those promises for you
as you type, before you ever run the program. Throughout this handbook you will see hints such as
\code{list[bool]} (a list of booleans), \code{tuple[int, int]} (a pair of integers --- one pixel
coordinate) and \code{list[dict]} (one dictionary per detected hand).

\begin{keyidea}
\textbf{Reading a signature aloud} is a genuinely useful skill. Take
\code{fingers\_state(pixel\_landmarks, hand\_label="Unknown", palm\_z=None, thumb\_mode="distance")}.
Read as: ``give me the pixel landmarks; optionally tell me which hand this is; give me five booleans
back --- one per finger, thumb first.'' Once you can read signatures, the architecture of the whole
project becomes legible.
\end{keyidea}

\subsection{Object-oriented programming, simplified}
\label{sec:oop}

Everything so far has been \emph{data} (values) and \emph{procedures} (functions) kept apart. Object
orientation glues them together.

\begin{pythonconcept}{Class, object, self --- the three words that block beginners}
A \textbf{class} is a blueprint. An \textbf{object} (or \emph{instance}) is a thing built from that
blueprint. \code{self} is how the object refers to itself.

Analogy: a class is the architectural drawing of a house; an object is an actual house. You can build
many houses from one drawing, and each has its own front door colour --- its own data --- while
sharing the same structure.

\begin{lstlisting}
class HandTracker:
    def __init__(self, max_hands=2):
        self.max_hands = max_hands     # store on THIS object
        self.hands_data = []           # each object gets its own empty list

    def hand_count(self):
        return len(self.hands_data)

tracker = HandTracker()                # build an object from the blueprint
print(tracker.hand_count())            # 0
\end{lstlisting}

Line by line:

\begin{itemize}
  \item \code{class HandTracker:} defines the blueprint. By convention, class names are capitalised.
  \item \code{\_\_init\_\_} is the \textbf{constructor}: it runs automatically when the object is
        built. The name has two underscores on each side --- in Python that signals ``special method
        that the language itself calls''.
  \item \code{self} is the first parameter of every method. Python passes the object in
        automatically, so \code{tracker.hand\_count()} is really
        \code{HandTracker.hand\_count(tracker)}. You never pass \code{self} yourself; you only
        declare it.
  \item \code{self.hands\_data = []} creates a piece of state \emph{belonging to this object}.
        Crucially, each object gets its own list. Two trackers would not share one hand cache.
\end{itemize}

Why does this project need a class at all? Because detection has \textbf{state}: the MediaPipe model
object itself, the smoothing history, the gesture debouncers. A bare function would have to be handed
all of that on every call. A class lets the state live next to the methods that use it --- which is
why the detector is \code{HandTracker}, and \code{find\_hands}, \code{find\_positions} and
\code{fingers\_up} are its methods.
\end{pythonconcept}

\subsection{Five special constructs used in this project}

\subsubsection*{1. \code{deque} with a maximum length}

\begin{codecard}
\begin{lstlisting}
from collections import deque
from hand_logic import GestureDebouncer

history = deque(maxlen=5)
history.append("Fist")        # ['Fist']
history.append("Peace")       # ['Fist', 'Peace']
# ... after five appends ...
history.append("Fist")        # the OLDEST entry is silently evicted
\end{lstlisting}
\end{codecard}

A \code{deque} (``deck'', double-ended queue) is a list optimised for adding and removing at both
ends. The \code{maxlen} argument makes it a \textbf{sliding window}: it never grows beyond five
entries, and it never needs to be trimmed by hand. That single feature is the entire implementation of
the gesture debouncer in \code{hand_logic.py}: keep the last five gestures, count how many agree.

\begin{pythonconcept}{Why not a list?}
You could write \code{history.append(g); if len(history) > 5: history.pop(0)} --- but
\code{pop(0)} on a list shifts every remaining element, which is $O(n)$. A \code{deque} is designed
for that operation and evicts in $O(1)$. It is also harder to forget. In a 30-frames-per-second
loop, forgetting to trim a list is a slow memory leak that only shows up during a long demo.
\end{pythonconcept}

\subsubsection*{2. List comprehensions}

\begin{codecard}
\begin{lstlisting}
# the long way
pixels = []
for lm in landmarks:
    pixels.append((int(lm.x * width), int(lm.y * height)))

# the comprehension
pixels = [(int(lm.x * width), int(lm.y * height)) for lm in landmarks]
\end{lstlisting}
\end{codecard}

Both blocks do the same thing. The second is the idiom you will see in this codebase wherever one
collection is transformed into another. Read a comprehension as: ``build this expression, for each
item in this collection.''

\subsubsection*{3. \code{if \_\_name\_\_ == "\_\_main\_\_":}}

\begin{codecard}
\begin{lstlisting}
def main():
    ...

if __name__ == "__main__":
    raise SystemExit(main())
\end{lstlisting}
\end{codecard}

\begin{pythonconcept}{The most-asked Python question, answered properly}
When Python runs a file \emph{directly} (\code{python hand\_tracker.py}), it sets a built-in
variable called \code{\_\_name\_\_} to the string \code{"\_\_main\_\_"} inside that file.

When the same file is \emph{imported} by another file (\code{import hand\_tracker}), Python sets
\code{\_\_name\_\_} to the module's own name, \code{"hand\_tracker"}.

So \code{if \_\_name\_\_ == "\_\_main\_\_":} means \textbf{``run this block only if I am the program
being started, not if I am being used as a library.''}

This project needs exactly that distinction for three reasons:

\begin{itemize}
  \item The test suite imports \code{hand\_tracker} to test its helper functions. If the capture
        loop ran at import time, \code{pytest} would try to open your webcam.
  \item \code{raise SystemExit(main())} converts the integer that \code{main()} returns into the
        process exit code, so CI can tell success (0) from failure (1).
  \item It keeps the top of the file importable for tooling without side effects.
\end{itemize}
\end{pythonconcept}

\subsubsection*{4. Exceptions and \code{finally}}

\begin{codecard}
\begin{lstlisting}
try:
    while True:
        ok, frame = capture.read()
        ...
finally:
    capture.release()          # ALWAYS runs: normal exit, error, or Ctrl+C
    cv2.destroyAllWindows()
\end{lstlisting}
\end{codecard}

An \textbf{exception} is Python's way of saying ``I cannot continue''. \code{try} marks a region
where you expect one; \code{finally} marks cleanup that must happen no matter how the region ends.
The tracker's capture loop uses \code{finally} to release the camera and close the window. Without
it, an unplugged webcam or a \code{Ctrl+C} would leave the camera device locked, and the next run
would fail with ``could not open a webcam'' until you rebooted --- a classic and infuriating bug.

\subsubsection*{5. The walrus-free guard clause}

\begin{codecard}
\begin{lstlisting}
def find_positions(self, frame, hand_index=0):
    if 0 <= hand_index < len(self.hands_data):
        return list(self.hands_data[hand_index]["px"])
    return []
\end{lstlisting}
\end{codecard}

Note the \emph{chained comparison} \code{0 <= i < len(...)}. Python evaluates it as
\code{(0 <= i) and (i < len(...))}, which is why a negative index --- which would otherwise mean
``count from the end'' --- is rejected. This two-line pattern (``check, return early; otherwise fall
through to the safe default'') is how the project removes an entire class of crashes:
``IndexError: list index out of range'' is impossible here.

\section{How Digital Images Work}
\label{sec:images}

\subsection{Pixels, resolution and channels}

An image is a grid of coloured dots called \textbf{pixels} (short for ``picture elements''). The
grid has a width and a height, measured in pixels; together they are the \textbf{resolution}.

\begin{mlconcept}{From photons to numbers}
A camera sensor measures, for each of its tiny light-sensitive cells, how much light arrived. The
webcam in this project is asked for $640 \times 480$ cells (that is 307{,}200 pixels). For each cell
we get three intensities: how much \textbf{red}, how much \textbf{green}, how much \textbf{blue} light
arrived.

Each intensity is stored as one byte --- an integer from 0 (none) to 255 (full). Three bytes per
pixel, 307{,}200 pixels: about 0.9 megabytes per frame, and at 30 frames per second roughly
27 MB/s flowing through your program. This is exactly why the project caches work per frame
instead of recomputing it (Chapter~\ref{ch:architecture}).

A colour is therefore an ordered triple. \code{(255, 0, 0)} is pure red in RGB order;
\code{(0, 0, 0)} is black; \code{(255, 255, 255)} is white.
\end{mlconcept}

\subsection{Why OpenCV thinks in BGR}

Here is a fact that costs beginners hours: OpenCV --- the library that reads your webcam --- stores
colour channels in the order \textbf{blue, green, red}, not red, green, blue. It is a historical
inheritance from the 1990s BGR camera hardware OpenCV was written against.

MediaPipe, on the other hand, was trained on ordinary RGB images.

\begin{edgetrap}{The BGR/RGB mismatch: silently wrong, not loudly broken}
If you feed a BGR image to an RGB model, nothing crashes. Every colour is simply permuted: red and
blue swap. Skin becomes a strange blue-ish tint, and the detector's accuracy quietly degrades --- or
drops to zero --- on some skin tones and lighting conditions.

It is the worst kind of bug: no exception, no error message, just worse results. This is why the
pipeline contains exactly one explicit conversion, at exactly one place:

\begin{lstlisting}
rgb = cv2.cvtColor(frame, cv2.COLOR_BGR2RGB)
results = self.hands.process(rgb)
\end{lstlisting}

Everything {\em before} that line is BGR (because it came from the camera and will be shown by
\code{cv2.imshow}); everything {\em after} is MediaPipe's business. Screen output stays in BGR,
because \code{cv2.imshow} expects BGR. The conversion is a one-way door for the model only.
\end{edgetrap}

\subsection{A frame is a three-dimensional array}

This is the single most useful mental model in the whole project.

\begin{mlconcept}{The flipbook and the cube}
Think of video as a \textbf{flipbook}: each page is one still image, and flipping fast enough
produces motion. A single frame is one page.

Now zoom into one page. It is a grid of pixels. Add the three colour values per pixel and the page
becomes a \textbf{three-dimensional block of numbers}:

\begin{itemize}
  \item \textbf{axis 0} --- the rows (height): $480$ of them
  \item \textbf{axis 1} --- the columns (width): $640$ of them
  \item \textbf{axis 2} --- the channels: $3$ of them
\end{itemize}

In NumPy --- Python's numerical array library --- a frame is described by its \code{shape}, and you
can inspect it at any moment:

\begin{lstlisting}
>>> frame.shape
(480, 640, 3)
>>> frame[0, 0]          # top-left pixel, all three channels
array([ 42,  42,  45], dtype=uint8)
>>> frame[240, 320]      # the pixel at the centre of the frame
array([118, 147, 201], dtype=uint8)
>>> frame[240, 320, 0]   # just the BLUE value of that pixel
118
\end{lstlisting}

Note the order: \code{frame[y, x]}, \emph{row first, column second}. This is a genuine trap, and it
is the reason the project writes \code{height, width = frame.shape[:2]} rather than the other way
around. MediaPipe works in the opposite convention --- it reports normalized $(x, y)$ --- so every
conversion in this codebase is a place where the two conventions meet.
\end{mlconcept}

\begin{keyidea}
\textbf{One sentence to memorise:} a frame is a \code{(height, width, 3)} block of bytes, indexed
\code{[row, column, channel]}, and a landmark is a normalised \code{(x, y)} pair indexed the other
way round. Most ``my dot is in the wrong place'' bugs are one of these two facts being forgotten.
\end{keyidea}

\section{Machine Learning Demystified}
\label{sec:ml}

\subsection{Rules plus data, versus data plus answers}

\begin{mlconcept}{The inversion that defines machine learning}
\textbf{Traditional programming.} You write the rules, you feed in data, you get answers out.

\begin{center}
\code{rules + data $\rightarrow$ answers}
\end{center}

Example: ``if the pixel's blue value is greater than its red value, call it skin.'' You wrote the
rule. The computer applies it.

\textbf{Machine learning.} You supply data and the correct answers for that data. The computer
\emph{derives} the rules.

\begin{center}
\code{data + answers $\rightarrow$ rules}
\end{center}

Example: you show the computer 30{,}000 photographs of hands together with the exact pixel location
of every knuckle, and it works out its own rules for finding knuckles.

You cannot write the rule for ``where is the pinky knuckle'' by hand. There is no threshold on red,
green or blue that finds it. That is precisely why machine learning is used for this half of the
project --- and why the other half (deciding whether a finger is up) is ordinary geometry that
\emph{can} be written by hand, and therefore \emph{should} be.
\end{mlconcept}

\subsection{What a neural network actually is}

\begin{mlconcept}{Weights, layers, and feature detectors}
Strip away the vocabulary and a neural network is a very long chain of multiply-and-add operations,
with a non-linear squashing function between the steps. That is all.

\begin{itemize}
  \item A \textbf{weight} is one number that gets multiplied by an input. A small network has
        thousands; a large one has billions.
  \item A \textbf{layer} is a group of weights applied together, producing a new set of numbers from
        the previous set.
  \item A \textbf{feature} is a pattern that a layer reacts to. Early layers react to edges and
        colour gradients. Middle layers react to combinations such as ``a curved edge above a
        straight one''. Late layers react to whole structures: ``this cluster of edges is a
        knuckle''.
\end{itemize}

An analogy that holds up better than most: think of a factory production line. Raw pixels go in at
one end. At each station, workers check for one specific simple thing and pass on a summary. By the
end of the line, the summary is abstract enough to answer the question --- ``the pinky knuckle is at
$(0.31, 0.62)$''.

The important consequence for you: \textbf{each layer sees only the previous layer's summary}, never
the raw image. That is what makes the whole thing fast enough to run on a CPU at 30 FPS.
\end{mlconcept}

\subsection{Training versus inference}

\begin{mlconcept}{We only ever run the trained model}
\textbf{Training} is the process of searching for the weights. It requires a large labelled dataset,
a great deal of computation --- usually GPUs for hours or days --- and the ability to measure error
and correct it. Training is what Google did once, in order to ship MediaPipe.

\textbf{Inference} is using a network whose weights are already fixed. Given an input, produce an
output. No learning happens. It is a one-directional calculation.

This project is \textbf{100\% inference}. The weights arrive inside the MediaPipe wheel as a
\code{.tflite} file --- a few megabytes of numbers --- and the library loads them into memory when
you write:

\begin{lstlisting}
self.hands = mp_hands.Hands(
    static_image_mode=mode,
    max_num_hands=max_hands,
    min_detection_confidence=det_conf,
    min_tracking_confidence=track_conf,
)
\end{lstlisting}

Those four keyword arguments are \emph{not} training settings. They choose how confident the model
must be before it reports something. Chapter~\ref{ch:engine} covers what each one does and why the
defaults are $0.7$ and $0.6$.
\end{mlconcept}

\subsection{Normalized coordinates, and the conversion to pixels}

This is the last prerequisite, and it is where the first real mathematics of the project appears.

MediaPipe does not report pixel positions. It reports \textbf{normalized coordinates}: numbers in the
range $[0.0, 1.0]$, where $0.0$ is the left (or top) edge of the image and $1.0$ is the right (or
bottom) edge.

\begin{mathdeepdive}{From normalized coordinates to pixels}
\textbf{Why normalized?} Because the model must not care about resolution. Whether the frame is
$640 \times 480$ or $1920 \times 1080$, a hand in the middle of the picture is at
$x_{\text{norm}} = 0.5$. Normalizing makes one trained model valid on every camera, forever. The
model predicts \emph{relative} position; the exact pixel is an application concern.

\textbf{The conversion.} If a landmark sits at normalized position $x_{\text{norm}}$ on an image
$W$ pixels wide, its pixel column is

$$x_{\text{px}} = x_{\text{norm}} \cdot W$$

and, since the image is $H$ pixels tall,

$$y_{\text{px}} = y_{\text{norm}} \cdot H$$

\textbf{Worked example.} MediaPipe reports the index fingertip at
$(x_{\text{norm}}, y_{\text{norm}}) = (0.5469,\ 0.4896)$ on the default $640 \times 480$ frame:

\begin{align*}
x_{\text{px}} &= 0.5469 \times 640 = 350.0 \\
y_{\text{px}} &= 0.4896 \times 480 = 235.0
\end{align*}

so the fingertip is drawn at pixel $(350, 235)$ --- measured from the \emph{top-left} corner.

\textbf{Why the code truncates instead of rounding.} The project writes

\begin{lstlisting}
return [(int(lm.x * width), int(lm.y * height)) for lm in landmarks]
\end{lstlisting}

and \code{int()} \emph{truncates} ($350.9 \rightarrow 350$) rather than rounds
($350.9 \rightarrow 351$). For drawing a filled circle of radius 10 pixels, a half-pixel bias is
invisible, and truncation is measurably faster. What matters is that the value is an integer at all:
OpenCV's drawing functions accept floats, but integer pixel coordinates make the cache smaller and
the arithmetic exact when the geometry chapter computes distances and cross products.
\end{mathdeepdive}

\begin{keyidea}
\textbf{Chapter summary.} Python is an interpreted language whose lists, tuples and dictionaries give
the project its data shapes. An image is a \code{(height, width, 3)} block of bytes in BGR order, and
MediaPipe speaks normalized RGB. Machine learning inverts the usual programming equation --- data
plus answers produces rules --- and this project uses a model that has already been trained, only
ever running inference. Landmarks arrive normalized and must be multiplied by the frame size to
become pixels.
\end{keyidea}

\chapter{The Under-the-Hood ML Engine}
\label{ch:engine}

Chapter~\ref{ch:foundations} established that this project only ever runs \emph{inference}: the
weights were trained once, by someone else, and shipped inside the MediaPipe wheel. This chapter
opens the box. It explains the two-stage architecture that makes real-time hand tracking possible on
a CPU, what the model actually outputs, and the three dependency pins that will break your
installation if you ignore them.

\section{The two-stage pipeline}

A naive design would run one enormous neural network over the entire $640 \times 480$ frame, every
frame, looking for 21 landmarks. That design would be far too slow: most of the image is background,
and re-searching the whole frame 30 times a second wastes almost all of the computation.

MediaPipe's solution is to split the problem into a cheap ``where'' and an expensive ``what''.

\begin{mlconcept}{Stage 1 --- BlazePalm: find the hand, coarsely}
\textbf{BlazePalm} is a \emph{single-shot detector}. Its job is not to find fingers; it is to answer
one question, roughly: \emph{is there a palm here, and where?}

It works by sliding a set of predefined boxes --- called \textbf{anchors} --- over a downsampled
version of the frame, and for each anchor predicting two things:

\begin{enumerate}
  \item a \textbf{score}: how likely is it that this box contains a palm?
  \item an \textbf{offset}: how should the box be nudged and resized to fit better?
\end{enumerate}

It predicts a palm rather than a whole hand on purpose. A palm is rigid: its shape and size vary
little between people and poses, and --- critically --- it looks similar whether the hand is open,
closed, or rotated. A fist and an open hand have nearly the same palm box. That robustness is what
makes the detector reliable enough to hand off to stage two.

The output of stage 1 is one or more rectangles (plus a rotation angle, for hands tilted in the
image plane) that are \emph{likely} to contain a palm.
\end{mlconcept}

\begin{mlconcept}{Stage 2 --- the landmark model: find the 21 points, precisely}
Stage 2 takes the box from stage 1, \textbf{crops the frame to that box}, resizes the crop to the
network's fixed input size, and runs a second, more expensive network whose only job is regression on
21 specific points.

Because the crop is small and tightly centred on the hand, the network sees the fingers at high
effective resolution --- a fingertip may occupy 30 pixels in the crop where it occupied 5 in the full
frame. That is how it can predict a knuckle to a fraction of a pixel.

The output is 21 landmarks, each with three coordinates $(x, y, z)$ normalised to $[0, 1]$ relative
to the \emph{full} frame --- the library maps the crop coordinates back for you. The $z$ value is
relative depth, roughly ``how far in front of or behind the wrist plane'', and it is what makes this
a 3D system: fingers pointing at the camera still separate correctly in depth.
\end{mlconcept}

\begin{center}
\begin{tikzpicture}[
  box/.style={draw=slateline, fill=slatebg, rounded corners=1.2mm, align=center,
              inner sep=3mm, font=\small\sffamily, text=navy},
  lbl/.style={font=\scriptsize\sffamily, text=slatetext},
  ar/.style={-{Stealth[length=2.2mm]}, draw=cyanx, line width=1pt}
]
  % Full frame
  \draw[draw=slateline, fill=white, line width=1pt] (0,0) rectangle (3.2,2.4);
  \node[lbl] at (1.6,-0.24) {full frame $640\times480$ (downsampled)};
  \node[lbl, text=cyanx] at (1.6, 2.05) {stage 1: BlazePalm};

  % palm box inside frame
  \draw[draw=amberx, line width=1.2pt, fill=ambersoft]
        (1.0,0.55) rectangle (2.15,1.85);
  \node[lbl, text=amberx] at (1.57, 0.32) {palm box};

  % arrow to crop
  \draw[ar] (3.32,1.2) -- (4.18,1.2) node[midway, above, lbl] {crop};

  % Crop box
  \draw[draw=slateline, fill=white, line width=1pt] (4.3,0.35) rectangle (6.1,2.05);
  \node[lbl] at (5.2,-0.24) {crop, resized};
  \node[lbl, text=tealx] at (5.2, 2.05) {stage 2};

  % hand inside the crop
  \begin{scope}[shift={(5.02,0.75)}, scale=0.30]
    \handcoordinates
    \handbones
    \foreach \i in {0,...,20} {\node[joint, minimum size=4pt] at (L\i) {};}
    \foreach \i in {4,12,16,20} {\node[jointtip, minimum size=5pt] at (L\i) {};}
  \end{scope}

  % arrow to output
  \draw[ar] (6.22,1.2) -- (7.08,1.2) node[midway, above, lbl] {regress};

  % Output box
  \node[box, text width=2.5cm] (out) at (8.5,1.2)
    {21 landmarks\\ $(x,y,z)$\\ normalised};
  \node[lbl, text=tealx] at (8.5, 2.05) {stage 2 output};
\end{tikzpicture}
\captionof{figure}{The two-stage detection pipeline. Stage 1 searches the whole frame for a palm box
--- cheap, because the image is downsampled and the target is rigid. Stage 2 runs a heavier network
only inside that box. The split is what buys 30 FPS on a CPU.}
\label{fig:twostage}
\end{center}

\begin{keyidea}
\textbf{The economic argument for two stages.} Suppose the full frame costs 100 units of compute and
a crop costs 10. Running the expensive network full-frame costs 100 per frame. Running a cheap
detector once and the expensive network on a crop costs roughly $5 + 10 = 15$ --- and when the hand
is moving smoothly, MediaPipe can even skip the detector entirely for several frames and simply
\emph{track} the box, because it also predicts the hand's position in the next frame. That is the
difference between 8 FPS and 30 FPS, and it is why the design works at all.
\end{keyidea}

\section{What the model reports: the 21 landmarks}

MediaPipe numbers the 21 landmarks in a fixed order. \textbf{These indices are the vocabulary of the
entire project} --- every rule in Chapter~\ref{ch:math} is expressed as arithmetic on them, so they
are worth learning properly.

\begin{figure}[h]
\centering
\begin{tikzpicture}[scale=1.35, every node/.style={font=\scriptsize\sffamily}]
  \handcoordinates
  \handbones
  \foreach \i in {0,...,20} {\node[joint] at (L\i) {};}
  \foreach \i in {4,8,12,16,20} {\node[jointtip] at (L\i) {};}

  % emphasised landmarks
  \node[draw=redx, circle, inner sep=1.6pt, line width=0.9pt] at (L8) {};

  % labels
  \node[right=1pt, text=navy, font=\scriptsize\sffamily\bfseries] at (L0)  {0 wrist};
  \node[right=1pt, text=navy, font=\scriptsize\sffamily\bfseries] at (L4)  {4 thumb tip};
  \node[left=1pt,  text=navy, font=\scriptsize\sffamily\bfseries] at (L5)  {5 index MCP};
  \node[above right=1pt, text=navy, font=\scriptsize\sffamily\bfseries] at (L8) {8 index tip};
  \node[left=1pt,  text=navy, font=\scriptsize\sffamily\bfseries] at (L17) {17 pinky MCP};

  \node[right=1pt, text=slatetext] at (L12) {12};
  \node[left=1pt,  text=slatetext] at (L16) {16};
  \node[left=1pt,  text=slatetext] at (L20) {20};
  \node[right=1pt, text=slatetext] at (L6)  {6};
  \node[left=1pt,  text=slatetext] at (L10) {10};
  \node[left=1pt,  text=slatetext] at (L14) {14};
  \node[left=1pt,  text=slatetext] at (L18) {18};

  % annotations
  \draw[decorate, decoration={brace, amplitude=3pt, mirror}, draw=tealx]
        (1.35,0.95) -- (1.35,3.55) node[midway, right=4pt, text=tealx, align=left]
        {four long\\fingers};
\end{tikzpicture}
\caption{The MediaPipe hand topology. The bold labels are the landmarks this project reasons about
directly: the wrist anchors the palm normal, the thumb tip and IP joint drive the thumb rule, and
indices 5 and 17 define the palm plane. Tip landmarks are drawn in amber; the index tip --- the one
carrying the smoothed marker --- is ringed in red.}
\label{fig:landmarks}
\end{figure}

\subsection{The landmark table}

The naming convention is systematic, and worth internalising: \textbf{MCP} is the knuckle
(metacarpophalangeal joint), \textbf{PIP} the first finger joint, \textbf{DIP} the second, and
\textbf{TIP} the end of the finger. \textbf{CMC} is the base of the thumb, at the wrist.

\begin{center}
\rowcolors{2}{slatebg}{white}
\begin{tabularx}{\textwidth}{@{}r l l X@{}}
\toprule
\rowcolor{white}
\textbf{Index} & \textbf{Name} & \textbf{Short} & \textbf{Role in this project} \\
\midrule
0  & Wrist              & WRIST      & Origin of both palm vectors; defines the palm plane \\
1  & Thumb CMC          & ---        & \textit{(not used directly)} \\
2  & Thumb MCP          & ---        & \textit{(not used directly)} \\
3  & Thumb IP           & THUMB\_IP  & The thumb rule's ``comparison'' point \\
4  & Thumb tip          & THUMB\_TIP & The thumb rule's ``extended?'' point \\
5  & Index MCP          & INDEX\_MCP & Second palm vector endpoint; index-finger base \\
6  & Index PIP          & ---        & The knuckle index 8 must beat vertically \\
7  & Index DIP          & ---        & \textit{(not used directly)} \\
8  & Index tip          & INDEX\_TIP & The accentuated, EMA-smoothed marker \\
9  & Middle MCP         & ---        & \textit{(not used directly)} \\
10 & Middle PIP         & ---        & The knuckle index 12 must beat vertically \\
11 & Middle DIP         & ---        & \textit{(not used directly)} \\
12 & Middle tip         & ---        & Middle finger extension test \\
13 & Ring MCP           & ---        & \textit{(not used directly)} \\
14 & Ring PIP           & ---        & The knuckle index 16 must beat vertically \\
15 & Ring DIP           & ---        & \textit{(not used directly)} \\
16 & Ring tip           & ---        & Ring finger extension test \\
17 & Pinky MCP          & PINKY\_MCP & The thumb rule's fixed reference point \\
18 & Pinky PIP          & ---        & The knuckle index 20 must beat vertically \\
19 & Pinky DIP          & ---        & \textit{(not used directly)} \\
20 & Pinky tip          & ---        & Pinky finger extension test \\
\bottomrule
\end{tabularx}
\end{center}

\begin{keyidea}
\textbf{Only nine of the twenty-one landmarks do any work.} The project reads indices
$0, 3, 4, 5, 8, 12, 16, 17, 20$ --- plus the four PIP joints $6, 10, 14, 18$ as vertical references.
That is a feature, not an oversight: a rule that uses fewer inputs is easier to reason about, easier
to test, and less exposed to jitter in landmarks that carry no information.
\end{keyidea}

\subsection{Why 21 points, and not a bounding box or a mask}

\begin{mlconcept}{The representation choice, and what it buys you}
There are three common ways to represent a detected hand:

\begin{itemize}
  \item \textbf{A bounding box} --- four numbers. Cheap, but useless for fingers: a fist and an open
        hand have the same box.
  \item \textbf{A segmentation mask} --- a pixel-per-pixel ``is this hand?'' image. Precise, but
        expensive, and it still does not tell you \emph{which} pixel is a knuckle.
  \item \textbf{Keypoints} --- a small, fixed set of semantically named points. Compact, cheap to
        compare, and every point has a \emph{meaning}: index 8 \emph{is} the index fingertip.
\end{itemize}

The project needs semantics, not pixels. Its questions are ``is this fingertip above that knuckle?''
and ``how far is the thumb from the pinky knuckle?'' --- questions about \emph{identified joints}.
Keypoints answer them directly with 63 numbers ($21 \times 3$) instead of a mask with 307{,}200.
\end{mlconcept}

\section{The dependency matrix, and three package traps}
\label{sec:deps}

The project pins exactly three packages:

\begin{center}
\rowcolors{2}{slatebg}{white}
\begin{tabularx}{\textwidth}{@{}l l X@{}}
\toprule
\rowcolor{white}
\textbf{Pin} & \textbf{Installs} & \textbf{Why exactly this} \\
\midrule
\code{mediapipe==0.10.14} & the model + runtime &
  The last release family with the legacy \code{mp.solutions.hands} API the project is built on.
  Everything from 0.10.30 onward has removed it. \\
\code{opencv-contrib-python}\newline\code{==4.11.0.86} & \code{cv2} &
  The \emph{contrib} build, not \code{opencv-python}: MediaPipe already depends on it, so pinning
  the same package avoids two competing \code{cv2} installs. \\
\code{numpy<2} & numerical arrays &
  MediaPipe 0.10.x wheels are compiled against the NumPy 1.x ABI. NumPy 2.0 changed that ABI. \\
\bottomrule
\end{tabularx}
\end{center}

Each of the three is a real trap that has cost real people real hours. Take them one at a time.

\begin{edgetrap}{Trap 1 --- two \code{cv2} packages fighting over one name}
OpenCV ships under several PyPI names: \code{opencv-python}, \code{opencv-contrib-python},
\code{opencv-python-headless}, and so on. They all install the \textbf{same importable module name},
\code{cv2}, into \code{site-packages}.

\textbf{MediaPipe's own dependency list includes \code{opencv-contrib-python}.} So this
\code{requirements.txt} would be actively harmful:

\begin{lstlisting}
mediapipe==0.10.14
opencv-python==4.10.0.84      # <-- WRONG: a second, competing cv2
numpy<2
\end{lstlisting}

pip will happily install both. Then \code{import cv2} resolves to whichever package wrote its files
last, and upgrading one can silently break the other --- producing the notorious
``it works on my machine'' failures that are almost impossible to debug because
\code{cv2.\_\_version\_\_} can report one thing while the code that runs comes from the other.

\textbf{The fix.} Pin the package MediaPipe already wants. The contrib build is a strict
\emph{superset} of the base build --- it adds extra modules --- so nothing is lost:

\begin{lstlisting}
mediapipe==0.10.14
opencv-contrib-python==4.11.0.86
numpy<2
\end{lstlisting}
\end{edgetrap}

\begin{edgetrap}{Trap 2 --- NumPy 2.0 and the ABI break}
NumPy 2.0 (released mid-2024) changed its C application binary interface. Extension modules compiled
against NumPy 1.x --- which includes the MediaPipe 0.10.x wheels --- can fail at import time with
messages about binary incompatibility, or worse, import and then crash.

Pinning \code{numpy<2} resolves it to the 1.26.x line, which is also the last version supporting
Python 3.9--3.12 across the board. If you see \code{pip} proposing to upgrade NumPy as a
``helpful'' side effect of another install, decline.
\end{edgetrap}

\begin{edgetrap}{Trap 3 --- \code{mp.solutions} no longer exists (the load-bearing pin)}
This is the one that matters most, because it is not a subtlety --- it is a hard removal.

The legacy API used by this project, \code{mp.solutions.hands.Hands}, was deprecated and then
\textbf{removed} in MediaPipe 0.10.30+. On a newer version the very first attribute access fails:

\begin{lstlisting}
>>> import mediapipe as mp
>>> mp.solutions.hands
AttributeError: module 'mediapipe' has no attribute 'solutions'
\end{lstlisting}

\textbf{Consequence:} \code{pip install -U mediapipe} bricks the application. The pin
\code{mediapipe==0.10.14} is therefore not a preference --- it is a load-bearing structural element,
in the same sense as a beam holding up a floor.

\textbf{Defence.} Rather than let the user discover this through an opaque
\code{AttributeError} three imports deep, the application checks at import time and fails loudly with
an actionable message. Here is the real guard, from \code{hand\_tracker.py}:

\begin{codecard}
\begin{lstlisting}
if not hasattr(mp, "solutions") or not hasattr(mp.solutions, "hands"):
    raise RuntimeError(
        "mediapipe>=0.10.30 removed mp.solutions, which this app is built on. "
        "Install the pinned version instead: pip install -r requirements.txt "
        "(do NOT `pip install -U mediapipe`)."
    )
\end{lstlisting}
\end{codecard}

\begin{itemize}
  \item \code{hasattr(mp, "solutions")} probes the attribute \emph{without} raising --- the whole
        point of \code{hasattr} is that it returns \code{False} instead of blowing up. This is
        exactly the situation it was designed for.
  \item The \code{or} short-circuits: if \code{solutions} is missing entirely, the second
        \code{hasattr} is never evaluated, so the guard itself cannot crash.
  \item \code{raise RuntimeError} rather than \code{sys.exit} keeps the failure inside Python's
        normal error machinery, so a test harness or an editor sees a real traceback.
  \item The message names \emph{the fix}, not just the problem. A good error message tells you what
        to type next.
\end{itemize}
\end{edgetrap}

\subsection{Suppressing the deprecation noise}

Deprecated APIs like these emit warnings on import, which pollute a live console that is already
printing FPS telemetry. The project silences them \emph{only around the import}, so genuine warnings
elsewhere still surface:

\begin{codecard}
\begin{lstlisting}
with warnings.catch_warnings():
    warnings.simplefilter("ignore")
    import mediapipe as mp
\end{lstlisting}
\end{codecard}

The \code{with} block restores the previous warning behaviour on exit. Warning filters are global
state; changing them permanently to hide one import's noise would suppress diagnostics you may
actually need later --- for instance a NumPy deprecation that only appears under load.

\section{Confidence thresholds: what 0.7 and 0.6 actually control}

The tracker constructs the model with four arguments. Two of them are thresholds, and they are the
difference between a twitchy detector and a sluggish one.

\begin{codecard}
\begin{lstlisting}
self.hands = mp_hands.Hands(
    static_image_mode=mode,
    max_num_hands=max_hands,
    min_detection_confidence=det_conf,     # default 0.7
    min_tracking_confidence=track_conf,    # default 0.6
)
\end{lstlisting}
\end{codecard}

\begin{mlconcept}{Two thresholds, two different questions}
\begin{itemize}
  \item \code{min\_detection\_confidence = 0.7} applies when the system is \emph{looking for a new
        hand} --- stage 1, the palm detector running over the whole frame. A candidate box must score
        at least 0.7 to be accepted.
  \item \code{min\_tracking\_confidence = 0.6} applies when the system is \emph{following a hand it
        already found}, using the landmarks predicted from the previous frame. The bar is lower
        because tracking has more evidence: it knows where the hand just was.
\end{itemize}

Why the asymmetry? Detection is a \emph{false-positive} problem: if you accept weak candidates, a
folded sleeve or a face can spawn a phantom hand. Tracking is a \emph{false-negative} problem: if you
drop the hand for one frame, the on-screen count flickers and the debouncer restarts.

So the design is deliberately conservative about discovering hands (0.7) and forgiving about keeping
them (0.6). Lowering detection to ~0.5 is the usual first experiment for a dim room --- with the
warning that phantom hands become likelier. Raising it toward 0.9 makes the tracker refuse to see
your hand until it is well lit and centred.
\end{mlconcept}

\subsection{\code{static\_image\_mode} and \code{max\_num\_hands}}

\begin{itemize}
  \item \code{static\_image\_mode=True} disables temporal tracking. Every frame is treated as an
        unrelated still, so the detector runs every time. It is slower and noisier on video, but it
        is exactly right for \code{--image} mode, where there is only one frame and nothing to
        track. The CLI exposes it as \code{--static} for live use; \code{run\_image()} always sets
        it to \code{True}.
  \item \code{max\_num\_hands=2} caps how many hands are searched for. It is a genuine performance
        knob: the detector's cost grows with the number of candidate boxes it keeps. Two is the
        right default for a two-handed demo; raise it only if you need more.
\end{itemize}

\section{Handedness, and why the label inverts}

The model also emits a \textbf{handedness} classification --- ``Left'' or ``Right'' --- with a score.
Beginners assume this is a property of the hand in the image. It is not: it is a prediction made
under an assumption about how the image was captured.

\begin{edgetrap}{The mirrored-feed trap: MediaPipe assumes a selfie}
MediaPipe's documentation is explicit: handedness is predicted \emph{assuming the input image is
mirrored}, as when a front-facing camera feed has already been flipped horizontally --- the normal
convention for a selfie.

This project flips the frame for exactly that reason --- a mirrored view is what makes a webcam feel
like a mirror and makes ``move right'' move right on screen:

\begin{lstlisting}
frame = cv2.flip(frame, 1)
\end{lstlisting}

But that flip happens \textbf{before} MediaPipe ever sees the frame. The result: on many webcams and
configurations the label comes back \emph{inverted} --- your physical right hand is reported as
``Left''.

\textbf{Why this is handled by design rather than patched.} The project never lets the label drive
geometry. The finger logic uses handedness only in the explicitly opt-in \code{--thumb-mode x}; the
default thumb rule is handedness-free by construction (Section~\ref{sec:thumb}). The label is
treated as \emph{display telemetry}, and the escape hatch is a display-only fix:

\begin{lstlisting}
--swap-handedness     # swap the Left/Right labels in the HUD
\end{lstlisting}

Because the swap happens inside \code{parse\_hand\_data}, it also flows through to the per-hand
smoothing and debouncing keys, so no state is orphaned by the change.
\end{edgetrap}

\begin{keyidea}
\textbf{Chapter summary.} Detection is deliberately split in two: a cheap detector finds a palm box,
and an expensive network regresses 21 named 3D landmarks inside that crop --- which is what makes
real-time performance on a CPU achievable. The three pins each defend against a specific, real
failure: two competing \code{cv2} packages, the NumPy 2 ABI break, and the removal of
\code{mp.solutions}. The two confidence thresholds are tuned asymmetrically --- strict about
discovering hands, forgiving about losing them --- and the handedness label is display-only, because
on a mirrored feed it cannot be trusted.
\end{keyidea}

\chapter{Mathematical and Geometric Algorithms}
\label{ch:math}

Everything the camera sees has now become 21 points with names (Chapter~\ref{ch:engine}). This
chapter is the part of the project that is \emph{ours}: turning those points into answers. It is the
only chapter where the reasoning is original rather than borrowed from a library, and it is where
every interesting bug lives.

The chapter moves from the simplest question to the most subtle:

\begin{enumerate}
  \item Which way is ``up'' on a screen, and why does that invert the finger test?
        (\S\ref{sec:screencs})
  \item Is the thumb out? --- the rule that survives rotation, and the physical pose that still
        defeats it. (\S\ref{sec:thumb})
  \item Which side of the hand is facing the camera? (\S\ref{sec:palm})
  \item How do we stop the numbers shaking? (\S\ref{sec:ema})
  \item What gesture is this, and how do we stop the label flickering? (\S\ref{sec:gesture})
\end{enumerate}

\section{The screen coordinate system, and the inverted y-axis}
\label{sec:screencs}

\begin{mathdeepdive}{Two coordinate systems, one picture}
In the mathematics you learned at school, the \textbf{Cartesian plane} has its origin in the middle
or at the bottom-left, the $x$-axis grows to the right, and \textbf{$y$ grows upward}. A point
``higher up'' has a \emph{larger} $y$.

A computer screen inherited a different convention from the cathode-ray tubes of the 1960s: the
electron beam scanned from the top-left corner, so that became the origin, and the beam's downward
sweep became the positive $y$ direction. **On a screen, $y$ grows downward.** A point higher up has a
\emph{smaller} $y$.

The two systems are related by a vertical flip:

$$y_{\text{cartesian}} = H - y_{\text{screen}}$$

where $H$ is the image height in pixels. Nothing else changes; $x$ grows to the right in both.
\end{mathdeepdive}

\begin{figure}[h]
\centering
\begin{tikzpicture}[every node/.style={font=\small\sffamily}]
  % ---- Cartesian (left) ----
  \begin{scope}[shift={(0,0)}]
    \draw[slateline, line width=0.8pt] (-0.4,-0.4) rectangle (3.6,3.0);
    \draw[-{Stealth[length=2mm]}, navy, line width=1pt] (0,0) -- (3.3,0)
          node[right, text=navy] {$x$};
    \draw[-{Stealth[length=2mm]}, navy, line width=1pt] (0,0) -- (0,2.7)
          node[above, text=navy] {$y$};
    \fill[redx] (2.0,1.8) circle (2.2pt);
    \node[right=2pt, text=redx] at (2.0,1.8) {$P$};
    \draw[cyanx, dashed] (0,1.8) -- (2.0,1.8);
    \draw[cyanx, dashed] (2.0,0) -- (2.0,1.8);
    \node[below, text=cyanx] at (1.0,0.05) {$x$};
    \node[left, text=cyanx] at (0.05,0.9) {$y$};
    \node[below=0.35cm, text=navy, font=\small\sffamily\bfseries] at (1.6,-0.5)
         {Cartesian: $y$ grows \emph{up}};
  \end{scope}

  % ---- Screen (right) ----
  \begin{scope}[shift={(5.4,0)}]
    \draw[slateline, line width=0.8pt, fill=slatebg] (-0.4,-0.4) rectangle (3.6,3.0);
    \draw[-{Stealth[length=2mm]}, navy, line width=1pt] (-0.4,-0.4) -- (3.3,-0.4)
          node[right, text=navy] {$x$};
    \draw[-{Stealth[length=2mm]}, navy, line width=1pt] (-0.4,-0.4) -- (-0.4,2.7);
    \node[above, text=navy] at (-0.4,2.7) {$y$};
    \node[below left, text=navy, font=\scriptsize] at (-0.4,-0.4) {$(0,0)$};
    \fill[redx] (2.0,1.3) circle (2.2pt);
    \node[right=2pt, text=redx] at (2.0,1.3) {$P$};
    \draw[cyanx, dashed] (-0.4,1.3) -- (2.0,1.3);
    \draw[cyanx, dashed] (2.0,-0.4) -- (2.0,1.3);
    \node[left, text=cyanx] at (0.75,0.85) {$y$};
    \node[below, text=cyanx] at (2.0,-0.15) {$\dots$};
    \node[below=0.35cm, text=navy, font=\small\sffamily\bfseries] at (1.6,-0.5)
         {Screen: $y$ grows \emph{down}};
    \node[right=0.1cm, text=slatetext, font=\scriptsize, align=left] at (3.4,1.3)
         {origin\\top-left};
  \end{scope}

  % relationship
  \draw[-{Stealth[length=2mm]}, tealx, line width=1pt] (4.0,1.3) -- (4.9,1.3);
  \node[text=tealx, font=\small\sffamily, above] at (4.45,1.35) {$H-y$};
\end{tikzpicture}
\caption{The same point $P$ in the two conventions. In the Cartesian frame the vertical arrow
points up and $P$ has a positive $y$. On a screen the vertical arrow points down and the identical
point now has $y$ measured from the top-left corner. Every ``my marker is mirrored vertically'' bug
is this figure being forgotten.}
\label{fig:coords}
\end{figure}

\subsection{Why a fingertip being ``above'' its knuckle means a smaller $y$}

This single fact is the whole basis of the four-finger test.

\begin{mathdeepdive}{Deriving the finger test}
Hold your hand up in front of the camera with the index finger extended. In the image, the
fingertip is physically higher than the knuckle.

Because $y$ grows \emph{downward}, ``physically higher'' means ``numerically smaller'':

$$y_{\text{tip}} < y_{\text{knuckle}} \quad \Longleftrightarrow \quad \text{the fingertip is higher on screen}$$

Now bend the finger so the tip curls down toward the palm. The tip drops below the knuckle in the
image, so its $y$ becomes \emph{larger}:

$$y_{\text{tip}} > y_{\text{knuckle}} \quad \Longleftrightarrow \quad \text{the finger is folded}$$

The test is therefore:

$$\text{finger extended} \iff y_{\text{tip}} < y_{\text{pip}}$$

and the comparison is \emph{strict} --- an exactly equal $y$ counts as folded, which costs nothing
and keeps the rule total (it always returns a definite answer).

\textbf{Worked example.} Say the index fingertip is at pixel $(346, 235)$ and the index PIP joint at
$(350, 280)$. Then $y_{\text{tip}} = 235 < y_{\text{pip}} = 280$, so the finger is extended --- and
indeed $45$ pixels of difference is a clearly raised finger. If instead the tip were at
$(352, 305)$, then $305 > 280$ and the finger is folded.
\end{mathdeepdive}

The four long fingers are tested with four (tip, knuckle) index pairs from \code{hand_logic.py}:

\begin{codecard}
\begin{lstlisting}
FINGER_PAIRS = [(8, 6), (12, 10), (16, 14), (20, 18)]
\end{lstlisting}
\end{codecard}

\begin{center}
\rowcolors{2}{slatebg}{white}
\begin{tabular}{@{}l l l l@{}}
\toprule
\rowcolor{white}
\textbf{Finger} & \textbf{Tip index} & \textbf{Knuckle (PIP) index} & \textbf{Test} \\
\midrule
Index  & 8  & 6  & $y_8 < y_6$ \\
Middle & 12 & 10 & $y_{12} < y_{10}$ \\
Ring   & 16 & 14 & $y_{16} < y_{14}$ \\
Pinky  & 20 & 18 & $y_{20} < y_{18}$ \\
\bottomrule
\end{tabular}
\end{center}

\begin{edgetrap}{The y-rule's hidden assumption: fingers must point roughly up}
The test compares $y$ only. That is a comparison of \emph{vertical} position, and it is only
meaningful while the fingers point roughly toward the top of the image.

Point your index finger horizontally to the right (thumb up, like a finger-gun) and the tip and
knuckle have almost the same $y$. The tip may even be slightly lower. The rule will report the finger
as folded while it is plainly extended.

This is a genuine limitation, not a bug, and it is why the tracker is documented as a
\emph{upright-hand} tool. Making it fully rotation-invariant would require comparing the fingertip
against the palm axis --- computing the finger's direction relative to the wrist-to-middle-knuckle
vector and asking whether the tip is farther along that axis than the knuckle. That is a real
technique, and the thumb rule below uses exactly that family of idea (a distance from a fixed palm
reference), but applying it to all four fingers is more machinery than this project needs.

\textbf{Practical consequence for you as a user:} keep your hand upright. \textbf{Practical
consequence as an engineer:} if you extend this project, this is the first rule to replace.
\end{edgetrap}

\section{The thumb, and the search for invariance}
\label{sec:thumb}

The thumb is the hard one. It moves on a different axis from the other four fingers, it is
foreshortened by perspective when it points at the camera, and its apparent left-right position
depends on which hand it is \emph{and} which way the palm is facing.

\subsection{The obvious rule, and exactly how it breaks}

The original specification proposed a horizontal comparison. Because a webcam feed is mirrored, and
because MediaPipe's handedness label is predicted under a mirrored-feed assumption, the rule was
written as:

$$\text{Right} \Rightarrow \text{thumb extended if } x_{\text{tip}} < x_{\text{ip}}$$
$$\text{Left} \Rightarrow \text{thumb extended if } x_{\text{tip}} > x_{\text{ip}}$$

\begin{edgetrap}{Why the horizontal rule fails: palm inversion}
The rule silently encodes a physical assumption --- \emph{the palm faces the camera}. Turn your hand
around so the back faces the camera, and everything the rule knows becomes false.

Concretely: hold up your right hand, palm toward the camera, thumb out. The thumb points to one side
of the image. Now rotate your wrist 180\textdegree\ so the back of the hand faces the camera, keeping
the thumb extended. The thumb is now on the \emph{other} side of the image, while the hand is still
your right hand and the label still reads ``Right''. The comparison $x_{\text{tip}} < x_{\text{ip}}$
now returns the opposite answer.

The result: the finger count flickers as you rotate your hand, and the gesture classification
collapses with it. Rotation of the wrist alone --- no change in pose at all --- flips the answer.

There is a second, subtler failure: the comparison is purely horizontal, so a hand rotated in the
image plane (say, 45\textdegree) has a thumb whose $x_{\text{tip}} - x_{\text{ip}}$ difference is
scaled by $\cos 45^\circ \approx 0.71$, shrinking the margin and making the test fragile near the
decision boundary.
\end{edgetrap}

The rule is kept in the code because the specification asked for it, reachable as
\code{--thumb-mode x}, but it is not the default.

\subsection{The distance rule}

\begin{keyidea}
\textbf{The idea in one sentence.}
However the hand is rotated, an \textbf{extended thumb reaches away from the palm}, while a folded
thumb is tucked \emph{across} the palm --- so measure how far the thumb tip and the thumb joint sit
from a fixed palm reference, and see which is farther.
\end{keyidea}

The fixed reference is landmark 17, the pinky MCP --- the far corner of the palm, the point the thumb
moves \emph{toward} when it folds and \emph{away} from when it extends.

\begin{mathdeepdive}{The rule, derived and justified}
$$\text{thumb extended} \iff d\big(P_4,\ P_{17}\big) > d\big(P_3,\ P_{17}\big)$$

where $P_4$ is the thumb tip, $P_3$ is the thumb IP joint, $P_{17}$ is the pinky MCP, and $d$ is the
ordinary Euclidean distance

$$d(A,B) = \sqrt{(A_x - B_x)^2 + (A_y - B_y)^2}.$$

\textbf{Why a distance and not a coordinate.} A coordinate is tied to the frame of reference: $x$
means ``how far right on the screen'', which changes when the hand rotates. A distance is an
\emph{intrinsic} quantity --- it depends only on the two points themselves.

\textbf{The invariance argument, stated precisely.} Suppose the hand is rotated in the image plane by
an angle $\theta$ about any centre $C$. The rotation matrix is

$$R_\theta = \begin{pmatrix} \cos\theta & -\sin\theta \\ \sin\theta & \cos\theta \end{pmatrix},
\qquad R_\theta^{\mathsf T} R_\theta = I.$$

Every landmark is transformed as $P' = C + R_\theta (P - C)$. Therefore the difference
$P_4' - P_{17}' = R_\theta (P_4 - P_{17})$, and

\begin{align*}
d(P_4', P_{17}') &= \big\| R_\theta (P_4 - P_{17}) \big\| \\
                 &= \sqrt{(P_4-P_{17})^{\mathsf T} R_\theta^{\mathsf T} R_\theta (P_4-P_{17})} \\
                 &= \sqrt{(P_4-P_{17})^{\mathsf T} (P_4-P_{17})} = d(P_4, P_{17}).
\end{align*}

The distance is \emph{exactly} unchanged by any rotation, and the same holds for $d(P_3, P_{17})$.
Since both sides of the comparison are preserved, the comparison's result cannot change. The edge
case is equality --- the two distances are equal exactly when the thumb tip and IP joint lie on a
circle centred on the pinky knuckle, which happens only in a degenerate straight-thumb pose; the
strict \code{>} then reports ``folded'', which is the safe answer for a thumb that is not clearly
away from the palm.

\textbf{Handedness-independence.} Nothing in the expression refers to ``left'' or ``right''. It is a
statement about three points, so it cannot be affected by a label that may itself be wrong.
\end{mathdeepdive}

\begin{mathdeepdive}{Worked example, with numbers}
This is the actual synthetic hand from the project's test suite (a right hand, palm to camera,
fingers up, mirrored view).

\textbf{Thumb extended}
\begin{align*}
P_4 &= (430, 320) \quad \text{(thumb tip)} \\
P_3 &= (415, 345) \quad \text{(thumb IP)} \\
P_{17} &= (262, 335) \quad \text{(pinky MCP)}
\end{align*}

\begin{align*}
d(P_4, P_{17}) &= \sqrt{(430-262)^2 + (320-335)^2} \\
               &= \sqrt{168^2 + (-15)^2} = \sqrt{28224 + 225} = \sqrt{28449} \approx 168.7
\end{align*}

\begin{align*}
d(P_3, P_{17}) &= \sqrt{(415-262)^2 + (345-335)^2} \\
               &= \sqrt{153^2 + 10^2} = \sqrt{23409 + 100} = \sqrt{23509} \approx 153.3
\end{align*}

Since $168.7 > 153.3$, the thumb is reported \textbf{extended}. \checkmark

\textbf{Thumb folded.} Now the tip folds across the palm to $P_4 = (375, 350)$, with the IP joint at
$P_3 = (395, 355)$:

\begin{align*}
d(P_4, P_{17}) &= \sqrt{113^2 + 15^2} = \sqrt{12769+225} = \sqrt{12994} \approx 114.0 \\
d(P_3, P_{17}) &= \sqrt{133^2 + 20^2} = \sqrt{17689+400} = \sqrt{18089} \approx 134.5
\end{align*}

Now $114.0 < 134.5$, so the thumb is reported \textbf{folded}. \checkmark

Notice the sizes involved: the margins are roughly $15$ pixels and $20$ pixels respectively. That is
a comfortable margin at $640 \times 480$ --- large enough that landmark jitter of one or two pixels
cannot flip the answer.
\end{mathdeepdive}

\subsection{The knife hand: when the distance rule lies}

\begin{edgetrap}{The adducted-thumb false positive}
Hold your hand flat, fingers straight, thumb pressed tightly against the side of the index finger ---
the ``karate chop'' or ``knife hand''.

The thumb tip is physically at the far end of the hand, near the index fingertip, while the thumb IP
joint is back near the palm. Distance to the pinky knuckle (landmark 17): the tip really \emph{is}
farther away than the IP joint. So the rule reports the thumb as \textbf{extended}, and the finger
count is one too high.

This is not a coding error; it is the rule doing exactly what it says. The failure is that the rule
only ever asks ``is the thumb away from the pinky side of the palm?'', and in a knife hand the answer
is genuinely yes.

\textbf{The candidate fix} --- a secondary clearance gate against the index MCP (landmark 5), which
the thumb must also be distant from:

$$\text{dist}(P_4, P_5) > 1.2 \cdot \text{dist}(P_3, P_5)$$

The factor $1.2$ encodes ``the tip must be \emph{clearly} farther from the index knuckle than the IP
joint, not merely farther''. In a knife hand the tip sits close to the index finger, so this second
test fails and the false positive is suppressed.

\textbf{Why this project does not enable it by default.} The gate trades one error for another. On a
thumb pointing sideways across the image --- a genuine ``thumbs up'' rotated 90\textdegree\ by wrist
tilt --- the tip can pass close to the index knuckle and the gate would suppress a \emph{correct}
extended result. The factor $1.2$ is a tuned constant with no derivation, and it varies with hand
size and camera distance because it compares raw pixel distances.

\textbf{The engineering position taken here:} ship the rule that is rotation-invariant and
documented, document the pose that defeats it, and leave the tuning to whoever has the hardware in
front of them. The limitation is stated plainly in the project's README rather than hidden.
\end{edgetrap}

\subsection{Keeping the old rule reachable, and what happens with no label}

The code implements both rules, with the distance rule as the default and a defined behaviour when
the handedness label is unavailable:

\begin{codecard}
\begin{lstlisting}
if thumb_mode == THUMB_X and hand_label in ("Left", "Right"):
    tip, ip = pixel_landmarks[THUMB_TIP], pixel_landmarks[THUMB_IP]
    thumb = tip[0] < ip[0] if hand_label == "Right" else tip[0] > ip[0]
else:
    thumb = _distance(pixel_landmarks[THUMB_TIP], pixel_landmarks[PINKY_MCP]) > _distance(
        pixel_landmarks[THUMB_IP], pixel_landmarks[PINKY_MCP]
    )
\end{lstlisting}
\end{codecard}

\begin{itemize}
  \item The condition is a conjunction: the naive rule needs \emph{both} \code{thumb\_mode == "x"}
        \emph{and} a usable label.
  \item If the label is \code{"Unknown"} --- which happens when MediaPipe's handedness list is
        missing or short --- the \code{else} branch applies \textbf{the distance rule}, not
        \code{False}. A missing label degrades to a working rule rather than silently losing the
        thumb.
  \item \code{tip[0]} is the $x$ component of the tip tuple; \code{tip[1]} would be $y$. Keeping
        coordinate access as indexed tuples rather than named objects is what makes this expression
        short.
\end{itemize}

\section{The palm normal: which side of the hand am I seeing?}
\label{sec:palm}

Finger counting does not need this. It is included because it is exactly the kind of diagnostic
telemetry the specification asks for, and because it makes the thumb rule's invariance visible while
the program runs.

\begin{mathdeepdive}{Two vectors and a cross product}
Take two vectors from the wrist $P_0$:

\begin{align*}
\vec{v}_1 &= P_5 - P_0 \quad \text{(wrist} \to \text{index knuckle)} \\
\vec{v}_2 &= P_{17} - P_0 \quad \text{(wrist} \to \text{pinky knuckle)}
\end{align*}

In 3D, the palm's normal vector is their cross product $\vec{v}_1 \times \vec{v}_2$. Writing
$\vec{v}_1 = (a_x, a_y, a_z)$ and $\vec{v}_2 = (b_x, b_y, b_z)$:

$$\vec{v}_1 \times \vec{v}_2 =
\begin{pmatrix} a_y b_z - a_z b_y \\ a_z b_x - a_x b_z \\ a_x b_y - a_y b_x \end{pmatrix}$$

The project uses \textbf{only the $z$ component}, because the landmarks' $x$ and $y$ are pixel
coordinates while $z$ is relative depth --- mixing them would be meaningless. With
$a = (x_5 - x_0,\ y_5 - y_0)$ and $b = (x_{17} - x_0,\ y_{17} - y_0)$:

\begin{align*}
Z_{\text{normal}} &= a_x b_y - a_y b_x \\
 &= (x_5 - x_0)(y_{17} - y_0) - (y_5 - y_0)(x_{17} - x_0)
\end{align*}

This is the signed area (times two) of the triangle wrist--index-knuckle--pinky-knuckle. Its
\textbf{sign} records the rotational order of the two vectors --- which way you turn going from
$\vec{v}_1$ to $\vec{v}_2$ --- and the two sides of a hand produce opposite orders.

\textbf{Worked example.} Using the test suite's synthetic palm-facing hand
$P_0 = (320,420)$, $P_5 = (355,320)$, $P_{17} = (262,335)$:

\begin{align*}
Z &= (355-320)(335-420) - (320-420)(262-320) \\
  &= (35)(-85) - (-100)(-58) \\
  &= -2975 - 5800 = -8775
\end{align*}

A \emph{negative} $Z$ for a palm facing the camera. Mirror the same hand horizontally (which is
exactly what ``showing the back of your hand'' does to the knuckle ordering) and every $x$ difference
changes sign, so $Z$ becomes $+8775$ --- the sign flips, as it must.
\end{mathdeepdive}

\begin{figure}[h]
\centering
\begin{tikzpicture}[scale=1.25, every node/.style={font=\scriptsize\sffamily}]
  \handcoordinates
  \handbones

  % palm triangle
  \fill[tealx, opacity=0.12] (L0) -- (L5) -- (L17) -- cycle;
  \draw[tealx, dashed, line width=0.7pt] (L0) -- (L5) -- (L17) -- cycle;

  \foreach \i in {0,...,20} {\node[joint] at (L\i) {};}
  \foreach \i in {4,8,12,16,20} {\node[jointtip] at (L\i) {};}

  % vectors
  \draw[-{Stealth[length=2.6mm]}, cyanx, line width=1.6pt] (L0) -- (L5)
        node[midway, above left=-1pt, text=cyanx] {$\vec{v}_1$};
  \draw[-{Stealth[length=2.6mm]}, amberx, line width=1.6pt] (L0) -- (L17)
        node[midway, above right=-1pt, text=amberx] {$\vec{v}_2$};

  \node[right=0pt, text=navy, font=\scriptsize\sffamily\bfseries] at (L0) {$P_0$};
  \node[right=1pt, text=navy, font=\scriptsize\sffamily\bfseries] at (L5) {$P_5$};
  \node[left=1pt, text=navy, font=\scriptsize\sffamily\bfseries] at (L17) {$P_{17}$};

  % normal indicator: circle-with-dot is the convention for "out of the page"
  \draw[violetx, line width=1.1pt] (-2.10, 1.30) circle (0.20);
  \fill[violetx] (-2.10, 1.30) circle (0.06);
  \node[violetx, font=\scriptsize\sffamily, align=center] at (-2.10, 0.72)
       {$\vec{n}=\vec{v}_1\times\vec{v}_2$\\\emph{out of the screen}};
  \draw[violetx, line width=0.7pt, dashed] (-1.88, 1.30) -- (0.35, 1.30);

  % sign summary
  \node[draw=slateline, fill=slatebg, rounded corners=1mm, inner sep=3pt, align=left]
        at (1.3,-1.05)
        {$Z<0 \Rightarrow$ \textbf{Palm} \quad $Z>0 \Rightarrow$ \textbf{Back}};
\end{tikzpicture}
\caption{The palm normal. The wrist and the two outer knuckles define a plane; the sign of the
$z$-component of $\vec{v}_1 \times \vec{v}_2$ says whether that plane is presented palm-first or
back-first. Because the calculation uses only two differences, it is one subtraction cheaper than a
true 3D cross product and needs no depth data.}
\label{fig:palmnormal}
\end{figure}

\subsection{The sign convention, and how to remember it}

The sign-to-label mapping depends on the mirroring of the image, so the project states its convention
explicitly: \textbf{a mirrored (selfie) view, as produced by \code{cv2.flip(frame, 1)}}.

\begin{mathdeepdive}{Where the sign comes from}
Consider a right hand, palm to camera, fingers up, in a horizontally mirrored image.

\begin{itemize}
  \item In a mirrored view your right hand appears on the \emph{right} side of the picture, and its
        thumb --- pointing away from the palm --- appears on the \emph{far right} of the hand.
  \item Therefore the index knuckle $P_5$ lies to the \emph{right} of the wrist, and the pinky
        knuckle $P_{17}$ lies to the \emph{left}: $x_5 > x_0$ and $x_{17} < x_0$.
  \item Both knuckles are \emph{above} the wrist, and $y$ grows downward, so
        $y_5 < y_0$ and $y_{17} < y_0$.
\end{itemize}

Substituting these signs into the formula from the previous box,
$a_x b_y = (+)(-)$ and $a_y b_x = (-)(-)$, gives
$Z = (\text{negative}) - (\text{positive}) < 0$.

Now turn the hand over. Index and pinky knuckles swap sides, so $x_5 < x_0$ and $x_{17} > x_0$;
both $y$ relations are unchanged. Now $a_x b_y$ is negative and $a_y b_x$ is positive, giving
$Z > 0$.

Hence, for the mirrored convention: $Z<0 \Rightarrow$ \textbf{Palm facing}, $Z>0 \Rightarrow$
\textbf{Back facing}. The project encodes exactly this:

\begin{lstlisting}
def orientation_from_palm_z(palm_z):
    """Map the palm-normal sign to the HUD label ``Palm'' or ``Back''."""
    return "Palm" if palm_z < 0 else "Back"
\end{lstlisting}

The degenerate case $Z = 0$ (wrist, index knuckle and pinky knuckle collinear --- geometrically
impossible for a real hand, but reachable with synthetic or corrupted data) falls to \code{Back}
through the \code{else}. A test asserts this, so the behaviour is fixed rather than accidental.
\end{mathdeepdive}

\begin{edgetrap}{The convention is tied to the mirror, not to the hand}
Because the sign depends on the image being flipped, the label inverts in \code{--image} mode, where
a still photograph is deliberately \emph{not} mirrored. A photograph of a palm-first hand in
\code{--image} mode will report \code{Back}.

That is not a bug in the geometry; it is the geometry faithfully reporting an unmirrored image with a
mirrored-image convention. It is documented in the README, and it matters only for the diagnostic
label --- no finger decision depends on it.
\end{edgetrap}

\section{Smoothing: the exponential moving average}
\label{sec:ema}

\begin{mathdeepdive}{Sensor jitter, and why raw coordinates shake}
Even with a completely still hand, the reported landmark positions move by one to three pixels from
frame to frame. The causes are mundane and unavoidable:

\begin{itemize}
  \item camera sensor noise, especially in low light;
  \item micro-movements of your hand that you cannot consciously control;
  \item the model's own regression error, which varies slightly as the hand's appearance in the crop
        changes.
\end{itemize}

At 30 FPS, a dot drawn from raw coordinates therefore vibrates visibly --- it reads as ``cheap
program''. The fix is a filter: blend each new measurement with the previous smoothed value.

$$\mathbf{P}^{(t)}_{\text{smooth}} = \alpha\,\mathbf{P}^{(t)}_{\text{raw}}
                                  + (1-\alpha)\,\mathbf{P}^{(t-1)}_{\text{smooth}},
\qquad \alpha = 0.65$$

Applied independently to $x$ and $y$. The name is \textbf{exponential moving average} (EMA), or
single-pole low-pass filter.

\textbf{Why it works.} Expand the recurrence backwards:

\begin{align*}
\mathbf{P}^{(t)} &= \alpha \mathbf{R}^{(t)} + (1-\alpha)\Big[\alpha \mathbf{R}^{(t-1)}
                    + (1-\alpha)\mathbf{P}^{(t-2)}\Big] \\
                 &= \alpha \mathbf{R}^{(t)} + \alpha(1-\alpha)\mathbf{R}^{(t-1)}
                    + (1-\alpha)^2 \mathbf{P}^{(t-2)} \\
                 &= \alpha \sum_{k=0}^{n} (1-\alpha)^k \mathbf{R}^{(t-k)}
                    + (1-\alpha)^{n+1}\mathbf{P}^{(t-n-1)}
\end{align*}

The current position is a weighted sum of \emph{all} past measurements, with geometrically decaying
weights $\alpha, \alpha(1-\alpha), \alpha(1-\alpha)^2, \dots$: recent frames matter most, old frames
fade out but never vanish entirely.

\textbf{The noise-reduction price.} The weights sum to $1$ in the limit, so a constant true position
is reproduced exactly. Random noise, however, is averaged down: the effective averaging window has
length roughly $1/\alpha \approx 1.5$ frames of ``memory mass'' concentrated in recent frames, and
the residual jitter is reduced by roughly $\sqrt{\alpha/(2-\alpha)} \approx 0.69$ for
$\alpha = 0.65$.

\textbf{The responsiveness price.} If the true position jumps (you flick your finger), the displayed
value approaches it geometrically and never exactly arrives. After one frame it has covered
$\alpha = 65\%$ of the gap; after two, $1-(1-\alpha)^2 = 87.75\%$; after three, $95.7\%$. Visible
lag at 30 FPS is therefore on the order of one to two frames --- about 30--60 milliseconds, below the
threshold at which a user perceives the marker as ``slow''.
\end{mathdeepdive}

\begin{mathdeepdive}{Worked example, and the choice of $\alpha$}
Suppose the smoothed value is $P^{(t-1)} = 300$ and the new raw measurement is $R^{(t)} = 340$.

$$P^{(t)} = 0.65 \times 340 + 0.35 \times 300 = 221 + 105 = 326$$

So the dot moves $26$ pixels of the $40$-pixel gap --- most of the way, but not all of it. The
remaining $14$ pixels will be consumed over the following frames.

Why $\alpha = 0.65$ rather than something else?

\begin{itemize}
  \item $\alpha \to 1$ (say $0.95$): almost no smoothing. Jitter is preserved; the dot is a raw
        measurement with a hat on.
  \item $\alpha \to 0$ (say $0.1$): heavy smoothing. The dot becomes a lazy, laggy ghost that
        visibly trails the finger.
  \item $\alpha = 0.65$: covers $65\%$ of any gap per frame --- fast enough to feel instantaneous at
        video rates, slow enough to flatten single-frame noise.
\end{itemize}

\textbf{First sighting.} On the very first frame there is no previous value, so the recurrence has no
starting point. Blending against zero would make the dot animate in from the top-left corner, which
looks like a bug. Instead, the implementation treats ``no previous'' as a special case and passes the
measurement through unchanged:

\begin{lstlisting}
def ema(previous, current, alpha=0.65):
    if previous is None:
        return (current[0], current[1])
    return (
        alpha * current[0] + (1 - alpha) * previous[0],
        alpha * current[1] + (1 - alpha) * previous[1],
    )
\end{lstlisting}
\end{mathdeepdive}

\begin{edgetrap}{Smoothing state is keyed by hand, not by list index}
MediaPipe does not guarantee the order of its hand list between frames. If your right hand leaves the
frame, or the two hands cross, the hand at index 0 in one frame may be the hand at index 1 in the
next.

If the smoothing history were stored per \emph{index} --- \code{\_smooth\_tip[0]},
\code{\_smooth\_tip[1]} --- a reorder would blend the position of one physical hand with the history
of the other. The visible symptom is a sudden teleport of the marker, which is precisely the
behaviour the filter exists to prevent.

The project keys the state by \textbf{handedness label} instead, falling back to the index only when
the label is unavailable:

\begin{lstlisting}
@staticmethod
def _state_key(hand_index, label):
    return label if label in ("Left", "Right") else f"index:{hand_index}"
\end{lstlisting}

This is not perfect --- handedness itself can flicker --- but it is meaningfully more stable than a
raw index, and it costs nothing. The same key is reused for the gesture debouncer, so a reorder
cannot contaminate gesture history either.
\end{edgetrap}

\section{Gestures and the debouncer}
\label{sec:gesture}

\subsection{From five booleans to a name}

The finger logic produces five booleans in a fixed order --- \code{[thumb, index, middle, ring,
pinky]}. A gesture is simply a recognised pattern in those five bits.

\begin{center}
\rowcolors{2}{slatebg}{white}
\begin{tabular}{@{}l c c c c c l@{}}
\toprule
\rowcolor{white}
\textbf{Gesture} & \textbf{Thumb} & \textbf{Index} & \textbf{Middle} & \textbf{Ring} & \textbf{Pinky}
 & \textbf{Notes} \\
\midrule
Open Palm  & \checkmark & \checkmark & \checkmark & \checkmark & \checkmark & all five extended \\
Fist       &            &            &            &            &            & all five folded \\
Thumbs Up  & \checkmark &            &            &            &            & thumb alone \\
Peace      &            & \checkmark & \checkmark &            &            & the V \\
Rock On    & \checkmark & \checkmark &            &            & \checkmark & index, pinky, thumb \\
(anything else) &   ---   &    ---     &    ---     &    ---     &    ---     & reported as \code{---} \\
\bottomrule
\end{tabular}
\end{center}

The order of the checks in the code matters, and is written so that the more specific patterns are
tested first. Note that ``Open Palm'' (all five true) and ``Thumbs Up'' (only the thumb true) cannot
collide, but reading the code top to bottom shows the intent:

\begin{codecard}
\begin{lstlisting}
if thumb and index and middle and ring and pinky:
    return "Open Palm"
if not (thumb or index or middle or ring or pinky):
    return "Fist"
if thumb and not (index or middle or ring or pinky):
    return "Thumbs Up"
if not thumb and index and middle and not (ring or pinky):
    return "Peace"
if thumb and index and not (middle or ring) and pinky:
    return "Rock On"
return NO_GESTURE
\end{lstlisting}
\end{codecard}

The negated conditions (``not ring and not pinky'') are what make each pattern exact rather than a
prefix match. Without them, ``Peace'' would also match a hand with all five fingers up.

\subsection{Why a raw per-frame answer flickers}

\begin{mathdeepdive}{Temporal majority voting}
During a transition --- say from Fist to Peace --- the fingers do not move in perfect synchrony, and
for a few frames the pattern belongs to no gesture at all. Classifying each frame independently
produces a label that flashes through intermediate states. Worse, a single noisy frame in the middle
of a steady pose can produce a one-frame spike to a wrong label.

The fix is a \textbf{rolling majority vote} over a short history. Formally, with window length $n$
and required agreement $k$, at frame $t$:

$$\text{report } g \iff \sum_{i=t-n+1}^{t} \mathbb{1}\!\left[g_i = g\right] \ \ge\ k,
\qquad n = 5,\ k = 4$$

where $\mathbb{1}[\cdot]$ is the indicator function (1 if the condition holds, 0 otherwise). If no
label reaches $k$, the previously reported label is held.

The implementation is a \code{deque} with \code{maxlen=5}, which gives the sliding window for free,
and a \code{count} on the newest label:

\begin{lstlisting}
def update(self, gesture):
    self._history.append(gesture)
    if self._history.count(gesture) >= self.min_agree:
        self._stable = gesture
    return self._stable
\end{lstlisting}

\textbf{Worked trace.} The window is the last five entries; entries are listed oldest first.

\begin{center}
\rowcolors{2}{slatebg}{white}
\begin{tabular}{@{}r l l l@{}}
\toprule
\rowcolor{white}
\textbf{Frame} & \textbf{Input} & \textbf{Window after append} & \textbf{Returned} \\
\midrule
1 & Peace & P            & \code{---} \\
2 & Peace & P P          & \code{---} \\
3 & Peace & P P P        & \code{---} \\
4 & Peace & P P P P      & \textbf{Peace} \\
5 & Peace & P P P P P    & Peace \\
6 & Fist  & P P P P F    & Peace \\
7 & Fist  & P P P F F    & Peace \\
8 & Fist  & P P F F F    & Peace \\
9 & Fist  & P F F F F    & \textbf{Fist} \\
\bottomrule
\end{tabular}
\end{center}

Reading it carefully: a change of gesture takes exactly four frames to be accepted, and a two-frame
blip can never break through, because two frames cannot reach a count of four in a five-slot window.
That is the precise trade: $\approx 130$\,ms of latency at 30 FPS, in exchange for the complete
elimination of single- and double-frame flicker.
\end{mathdeepdive}

\begin{keyidea}
\textbf{Warm-up behaviour is defined, not accidental.} Before any majority has formed, the stable
label is the placeholder \code{---}, so the HUD shows no gesture for the first three frames of a
session rather than guessing. The alternative --- returning the first raw label immediately --- would
make the debouncer pointless on startup, since the first frame is exactly the one most likely to be
noisy.
\end{keyidea}

\section{The complete decision path}

Putting the chapter together, one frame produces this chain, entirely inside \code{hand\_logic.py}:

\begin{center}
\begin{tikzpicture}[
  node distance=0.32cm,
  st/.style={draw=slateline, fill=slatebg, rounded corners=1mm, align=center,
             inner sep=2.2mm, font=\scriptsize\sffamily, text=navy},
  ar/.style={-{Stealth[length=1.8mm]}, draw=tealx, line width=0.8pt}
]
  \node[st] (a) {normalized\\landmarks};
  \node[st, right=of a] (b) {pixel landmarks\\\code{to_pixel_landmarks}};
  \node[st, right=of b] (c) {palm normal\\\code{palm_normal_z}};
  \node[st, right=of c] (d) {5 booleans\\\code{fingers_state}};
  \node[st, right=of d] (e) {gesture name\\\code{classify_gesture}};
  \node[st, right=of e] (f) {stable label\\\code{GestureDebouncer}};

  \draw[ar] (a) -- (b);
  \draw[ar] (b) -- (c);
  \draw[ar] (c) -- (d);
  \draw[ar] (d) -- (e);
  \draw[ar] (e) -- (f);
\end{tikzpicture}
\captionof{figure}{The per-frame decision chain in the pure-logic module. Every stage is a plain
function of its inputs --- no global state, no camera, no I/O --- which is why the whole chain is
unit-testable with hand-written coordinates.}
\end{center}

\begin{keyidea}
\textbf{Chapter summary.} On a screen $y$ grows downward, so a fingertip is extended when
$y_{\text{tip}} < y_{\text{pip}}$ for the four long fingers. The thumb uses a distance comparison
against the pinky knuckle, which is provably invariant under image-plane rotation and independent of
handedness --- with the adducted ``knife hand'' documented as its known failure. The palm normal's
sign reveals which side of the hand faces the camera. An exponential moving average with
$\alpha = 0.65$ removes jitter at the cost of roughly a frame of latency, and a five-frame,
four-vote rolling majority removes gesture flicker at the cost of roughly four frames of latency.
\end{keyidea}

\chapter{System Design and Production Architecture}
\label{ch:architecture}

The algorithms are settled. This chapter is about everything that makes them survive contact with a
real machine: how the code is split, why there is a cache, what happens when the camera is missing,
how the whole thing is tested without hardware, and how a text overlay is drawn so that it stays
readable.

\section{Architectural separation}
\label{sec:split}

The project is two modules, and the boundary between them is the most important design decision in
the codebase.

\begin{figure}[h]
\centering
\begin{tikzpicture}[
  mod/.style={draw=navy, fill=slatebg, rounded corners=2mm, align=center,
              inner sep=4mm, font=\small\sffamily, text=navy, line width=1pt},
  itm/.style={font=\scriptsize\sffamily, text=slatetext, align=left},
  ar/.style={-{Stealth[length=3mm]}, draw=cyanx, line width=1.4pt},
  arred/.style={-{Stealth[length=3mm]}, draw=redx, line width=1.6pt, dashed}
]
  % hand_tracker
  \node[mod, text width=6.4cm] (t) at (0,0)
    {\textbf{hand\_tracker.py}\\[2pt]
     \rule{6cm}{0.4pt}\\[3pt]
     \begin{tabular}{@{}l@{}}
       camera I/O \textbullet\ \code{cv2.imshow}\\
       BGR$\to$RGB conversion\\
       MediaPipe model ownership\\
       HUD rendering \textbullet\ CLI \textbullet\ logging\\
       fallback chains \textbullet\ teardown
     \end{tabular}};

  % hand_logic
  \node[mod, text width=6.4cm, draw=tealx] (l) at (8.6,0)
    {\textbf{hand\_logic.py}\\[2pt]
     \rule{6cm}{0.4pt}\\[3pt]
     \begin{tabular}{@{}l@{}}
       landmark parsing \textbullet\ cache building\\
       finger \& thumb rules\\
       palm normal \textbullet\ bounding box\\
       EMA \textbullet\ gesture classification\\
       debouncing
     \end{tabular}};

  % the dependency arrow
  \draw[ar] (t.east) -- node[above, font=\scriptsize\sffamily, text=cyanx]{imports} (l.west);

  % what is forbidden
  \draw[arred] (l.south) -- ++(0,-1.5) node[below, font=\scriptsize\sffamily, text=redx, align=center]
       {no \code{import cv2}\\ no \code{import mediapipe}\\ no camera, no disk, no clock};

  % what is gained
  \node[font=\scriptsize\sffamily, text=tealx, align=center] at (8.6,-2.5)
       {$\Rightarrow$ testable in a bare Python environment};

  % hardware
  \node[mod, text width=3.4cm, draw=slatetext, fill=white] (hw) at (-5.6,0)
    {\textbf{hardware}\\[2pt]\rule{3cm}{0.4pt}\\[3pt]
     \begin{tabular}{@{}l@{}}
       webcam device\\
       display server
     \end{tabular}};
  \draw[ar] (hw.east) -- (t.west);
\end{tikzpicture}
\caption{The dependency rule. All hardware, all third-party libraries and all I/O live on the left.
The right-hand module is pure: it can be imported, exercised and exhaustively tested in a bare Python
interpreter with no camera, no display and no OpenCV installed. The dashed red arrow states what must
never appear in it.}
\label{fig:architecture}
\end{figure}

\subsection{Why the split exists}

The original specification asked for a single \code{hand\_tracker.py} containing everything, with
``unit-testable helpers''. Those two requirements contradict each other: a helper that lives in a
file which imports \code{cv2} and \code{mediapipe} cannot be tested without installing both, and
importing the module to test one geometry function drags in a 40-megabyte native library and, on a
headless machine, may fail outright.

Splitting resolves the contradiction cleanly:

\begin{itemize}
  \item \code{hand\_logic.py} imports \code{math} and \code{collections} only. Its functions take
        plain numbers and lists and return plain numbers and lists.
  \item \code{hand\_tracker.py} owns everything that touches the outside world.
\end{itemize}

\begin{keyidea}
\textbf{The rule, stated as a test you can apply to any new line of code:} \emph{would this line work
if there were no camera, no screen and no MediaPipe?} If yes, it belongs in \code{hand\_logic.py}. If
no, it belongs in \code{hand\_tracker.py}.

Everything the project knows about \emph{hands} is on the pure side. Everything it knows about
\emph{computers} is on the other.
\end{keyidea}

\subsection{What the split buys, concretely}

\begin{center}
\rowcolors{2}{slatebg}{white}
\begin{tabularx}{\textwidth}{@{}l X@{}}
\toprule
\rowcolor{white}
\textbf{Benefit} & \textbf{Consequence in this project} \\
\midrule
Fast tests & The logic suite runs in under a second, with no model loading and no hardware. \\
CI without a GPU & The continuous-integration job needs \code{pip install pytest} and nothing else
                  --- no 200 MB wheel, no display server. \\
Deterministic tests & Synthetic coordinates produce exactly predictable answers; a real webcam never
                      would. \\
Portability & The geometry is valid on any platform, at any resolution, for any camera. \\
Debuggability & When the count is wrong you can ask ``is the geometry wrong, or is the input wrong?''
                and answer it by feeding the geometry a known-good point set. \\
\bottomrule
\end{tabularx}
\end{center}

The test suite makes the payoff visible: the majority of the tests construct arrays of coordinates by
hand and assert on booleans --- no hardware anywhere in sight.

\section{Single-pass telemetry caching}
\label{sec:cache}

\subsection{The problem: repeated work in a hot loop}

A naive implementation converts normalized landmarks to pixels wherever it needs them. In one frame,
the code needs pixel coordinates for: the index-tip marker, all five finger tests (six landmarks),
the thumb rule (three landmarks), the palm normal (three landmarks), and the bounding box (all 21).
Recomputing inside each of those means iterating the landmark list five or more times per hand, per
frame, at 30 frames per second --- thousands of redundant multiplications per second.

That is not a catastrophe. It is waste, and waste in a real-time loop is what separates a program
that holds 30 FPS on a modest laptop from one that drops frames under load.

\subsection{The solution: one parse, one dictionary per hand}

\begin{codecard}
\begin{lstlisting}
hands.append(
    {
        "label": label,
        "norm": norm,
        "px": px,
        "bbox": bbox_from_pixels(px),
        "palm_z": z,
        "orientation": orientation_from_palm_z(z),
    }
)
\end{lstlisting}
\end{codecard}

Everything downstream reads from this structure. The cost of the conversion is paid exactly once per
hand per frame.

\begin{center}
\rowcolors{2}{slatebg}{white}
\begin{tabularx}{\textwidth}{@{}l l X@{}}
\toprule
\rowcolor{white}
\textbf{Key} & \textbf{Type} & \textbf{What it stores, and who consumes it} \\
\midrule
\code{label} & \code{str} & \code{"Left"}, \code{"Right"} or \code{"Unknown"}, after any
  \code{swap\_handedness}. Shown in the HUD; used as the state key for smoothing and debouncing. \\
\code{norm} & \code{list[tuple]} & The raw $(x, y, z)$ triples in $[0,1]$. Kept because it is the
  model's own output --- useful for debugging, and free to retain. \\
\code{px} & \code{list[tuple]} & Integer pixel landmarks. Consumed by every geometric rule and by
  the tip marker. \\
\code{bbox} & \code{tuple} & $(x_{\min}, y_{\min}, x_{\max}, y_{\max})$. The hand's screen extent;
  used for labelling and diagnostics. \\
\code{palm\_z} & \code{float} & The signed palm-normal value from Chapter~\ref{ch:math}. \\
\code{orientation} & \code{str} & \code{"Palm"} or \code{"Back"} --- the label derived once from
  \code{palm\_z}, so no consumer has to repeat the sign test. \\
\bottomrule
\end{tabularx}
\end{center}

\begin{edgetrap}{The cache has a lifetime, and it is one frame}
\code{find\_positions} is called with a frame argument, but it never reads that argument --- the
pixels were computed inside the preceding \code{find\_hands} call, using \emph{that} frame's
dimensions. The parameter exists only to preserve the signature the specification asked for.

The consequence is a real constraint: \textbf{the cached pixel coordinates are only valid for a frame
of the size they were computed from}. Reuse the tracker against a differently sized frame without
calling \code{find\_hands} first and the coordinates will be scaled wrongly --- the dot will be drawn
in the wrong place, with no error raised.

The rule to follow: one \code{find\_hands} per frame, then read. Never cache across frames, never mix
sizes. This is documented in the method's docstring precisely because the failure is silent.
\end{edgetrap}

\subsection{Invalid input, and the safe result}

Parsing is defensive because the model's output shape can vary: a hand with fewer than 21 landmarks
cannot be reasoned about, and a missing handedness list is normal in some configurations.

\begin{codecard}
\begin{lstlisting}
hand_landmarks = getattr(results, "multi_hand_landmarks", None)
if not hand_landmarks:
    return []

handedness = getattr(results, "multi_handedness", None) or []
...
landmarks = getattr(hand, "landmark", None)
if not landmarks or len(landmarks) < EXPECTED_LANDMARKS:
    continue
\end{lstlisting}
\end{codecard}

Three defensive decisions, each with a reason:

\begin{itemize}
  \item \code{getattr(..., None)} rather than \code{results.multi\_hand\_landmarks}: a result object
        with no hands may legitimately lack the attribute. \code{getattr} returns the default instead
        of raising \code{AttributeError}.
  \item \code{... or []} handles the ``attribute exists but is \code{None}'' case, which would
        otherwise break \code{len()}.
  \item \code{continue} rather than \code{break} for a short landmark list: one malformed hand does
        not discard the other, valid hand in the same frame.
\end{itemize}

The same defensiveness applies to the finger rules, which return \code{[False] * 5} --- never raise,
never return \code{None} --- when handed nothing usable. That contract is asserted by a test, so
callers may rely on it.

\section{Camera fallback and hardware resiliency}
\label{sec:camera}

Opening a webcam is the least reliable operation in the entire program. The device may be absent, in
use by a video call, or present but only on a different index. A single-shot
\code{cv2.VideoCapture(0)} fails on all three.

\begin{figure}[h]
\centering
\begin{tikzpicture}[
  n/.style={draw=slateline, fill=slatebg, rounded corners=1.2mm, align=center,
            inner sep=2.4mm, font=\scriptsize\sffamily, text=navy, text width=2.7cm},
  a/.style={-{Stealth[length=2mm]}, draw=cyanx, line width=0.9pt},
  f/.style={-{Stealth[length=2mm]}, draw=redx, line width=0.9pt},
  y/.style={-{Stealth[length=2mm]}, draw=tealx, line width=1.2pt}
]
  \node[n] (start) {requested index\\\code{--camera N}};
  \node[n, right=0.9cm of start] (be) {for each backend:\\DSHOW $\to$ MSMF $\to$ ANY};
  \node[n, below=0.7cm of be] (next) {next candidate index\\(1, then 0)};
  \node[n, right=0.9cm of be] (ok) {\textbf{opened}\\set 640$\times$480\\log the index};
  \node[n, below=0.7cm of ok, draw=redx, fill=ambersoft] (fail)
      {\textbf{all candidates failed}\\log an actionable error\\exit code 1};

  \draw[a] (start) -- (be);
  \draw[a] (be) -- node[above, font=\scriptsize\sffamily, text=tealx]{isOpened} (ok);
  \draw[f] (be) -- node[right, font=\scriptsize\sffamily, text=redx, align=left]{not opened} (next);
  \draw[a] (next) -| (be);
  \draw[f] (next) -- (fail);
\end{tikzpicture}
\caption{The camera fallback chain. Both axes are searched --- backends within one index, and then
other indices --- before giving up with a message that names the likely cause.}
\label{fig:camera}
\end{figure}

\begin{codecard}
\begin{lstlisting}
for candidate in [index] + [other for other in (1, 0) if other != index]:
    for backend in backends:
        capture = cv2.VideoCapture(candidate, backend)
        if capture.isOpened():
            capture.set(cv2.CAP_PROP_FRAME_WIDTH, width)
            capture.set(cv2.CAP_PROP_FRAME_HEIGHT, height)
            LOG.info("Camera %d opened at %dx%d.", ...)
            return capture
        capture.release()
    LOG.warning("Camera %d did not open.", candidate)

LOG.error("Could not open a webcam. Check it is connected and not in use by another app.")
return None
\end{lstlisting}
\end{codecard}

\begin{itemize}
  \item \textbf{Backends differ per platform.} On Windows the list is
        \code{[CAP\_DSHOW, CAP\_MSMF, CAP\_ANY]}: DirectShow is the most permissive for consumer
        webcams, Media Foundation is the modern stack, and \code{CAP\_ANY} lets OpenCV choose.
        Elsewhere the list is just \code{[CAP\_ANY]}, because the alternatives do not exist.
  \item \textbf{The candidate list is} \code{[index] + [1, 0] with the requested index removed},
        so \code{--camera 0} tries 0 then 1, while \code{--camera 5} tries 5, then 1, then 0.
  \item \code{capture.release()} on every failure is not optional. Without it, each failed attempt
        leaks a device handle, and on Windows a leaked handle can itself prevent the next attempt
        from succeeding --- the fallback chain would sabotage itself.
  \item \textbf{Width and height are requested, not guaranteed.} \code{cap.set} is a request to the
        driver; the log line reports what was actually granted, which is why it re-reads the
        properties rather than echoing the requested numbers.
\end{itemize}

\begin{edgetrap}{``Could not open a webcam'' --- the four real causes}
The error text is deliberately specific because the causes are mundane and diagnosable:

\begin{enumerate}
  \item The camera is in use by another application (a video call, a browser tab holding the device,
        or a previous run of this program that was killed without cleanup).
  \item The index is wrong --- common on laptops with both an integrated and an external camera, and
        on desktops where a capture card occupies index 0.
  \item Operating-system privacy settings are blocking camera access for desktop applications.
  \item The device is genuinely absent or disabled in Device Manager.
\end{enumerate}

Cause 1 is the most common by a wide margin, and it is the reason teardown matters
(Section~\ref{sec:teardown}).
\end{edgetrap}

\section{Headless mode: verifying without a webcam}
\label{sec:headless}

Every test in the project's logic suite passes without a camera --- but until this feature existed,
\emph{the drawing and rendering path} could only be exercised by a human sitting in front of a
webcam. That is a large untested surface: the mesh drawing, the tip marker, the HUD blending, the
CLI, the image encoding.

\code{--image PATH} closes that gap. It runs the entire pipeline --- load, detect, parse, compute,
draw, write --- on a single still.

\begin{center}
\rowcolors{2}{slatebg}{white}
\begin{tabularx}{\textwidth}{@{}r l X@{}}
\toprule
\rowcolor{white}
\textbf{Step} & \textbf{Call} & \textbf{Purpose} \\
\midrule
1 & \code{cv2.imread(path)} & Returns \code{None} on failure; the code checks and exits non-zero. \\
2 & \code{cv2.resize} & Only if \code{--width} \emph{and} \code{--height} were supplied. A still has
    a native resolution worth preserving. \\
3 & \code{HandTracker(mode=True)} & \code{static\_image\_mode} --- there is no next frame to track
    into. \\
4 & \code{find\_hands} $\to$ \code{annotate\_hands} $\to$ \code{draw\_hud} & \emph{The same
    functions the live loop calls.} No parallel implementation exists to drift out of sync. \\
5 & \code{cv2.imwrite} & Writes \code{<stem>\_out.png} next to the input, then exits 0. \\
\bottomrule
\end{tabularx}
\end{center}

Note the deliberate omission: the still is \textbf{not mirrored}. A photograph is not a selfie, and
flipping it would move every landmark the model found to the wrong side of the image.

\subsection{Continuous integration without a display}

\begin{codecard}
\begin{lstlisting}[style=codetextonly]
name: CI
on: [push, pull_request]

jobs:
  logic-tests:
    runs-on: ubuntu-latest
    strategy:
      matrix:
        python-version: ["3.11", "3.12"]
    steps:
      - uses: actions/checkout@v4
      - uses: actions/setup-python@v5
        with: { python-version: "${{ matrix.python-version }}" }
      - run: pip install pytest
      - run: pytest -q

  pipeline-smoke:
    runs-on: ubuntu-latest
    steps:
      - uses: actions/checkout@v4
      - uses: actions/setup-python@v5
        with: { python-version: "3.12" }
      - run: sudo apt-get update && sudo apt-get install -y libgl1 libglib2.0-0
      - run: pip install -r requirements.txt pytest
      - run: pytest -q
      - run: python hand_tracker.py --image tests/fixtures/hand.png
      - run: test -f tests/fixtures/hand_out.png
\end{lstlisting}
\end{codecard}

\begin{edgetrap}{\code{libGL.so.1}: the OpenCV-on-CI gotcha}
The \code{pipeline-smoke} job installs the real dependencies --- including
\code{opencv-contrib-python}, which is \emph{not} the headless build. Those wheels link against
system graphics libraries that a minimal container does not have, and the failure is a stark

\begin{lstlisting}[style=codetextonly]
ImportError: libGL.so.1: cannot open shared object file: No such file or directory
\end{lstlisting}

It looks like a broken install or a code bug. It is neither: it is a missing shared library in the
runner image. Installing \code{libgl1} (and \code{libglib2.0-0}, which the same wheels frequently
need) fixes it in one line. This is the single most commonly hit OpenCV/CI failure, which is why it
is commented in the workflow itself rather than left as folklore.
\end{edgetrap}

\section{The alpha-blended HUD}
\label{sec:hud}

Drawing white text directly onto a webcam image fails the moment anything bright is behind it --- a
window, a lamp, a white wall. The fix is to darken a rectangle behind the text before writing it.

\begin{mathdeepdive}{Alpha blending}
For each pixel in the card region, the output is a weighted average of what was there and the
overlay:

$$I_{\text{out}} = (1-\alpha)\, I_{\text{roi}} + \alpha\, I_{\text{overlay}}$$

The project uses $\alpha = 0.6$ with a black overlay:

$$I_{\text{out}} = 0.4\, I_{\text{roi}} + 0.6 \times 0 = 0.4\, I_{\text{roi}}$$

So a bright background pixel of 255 becomes $102$ --- comfortably dark --- while a dark background
of 40 becomes $16$. The result is that the card is visible as a subtle darkening on any background,
and white text on top is readable at any luminance. Note that the formula is applied
\emph{per channel}, so hues are preserved and only brightness is scaled.
\end{mathdeepdive}

\begin{codecard}
\begin{lstlisting}
card = frame[y0:y1, x0:x1]
frame[y0:y1, x0:x1] = cv2.addWeighted(card, 1.0 - HUD_ALPHA, np.zeros_like(card), HUD_ALPHA, 0)
\end{lstlisting}
\end{codecard}

Two details worth noticing:

\begin{itemize}
  \item \code{frame[y0:y1, x0:x1]} is a \textbf{view}, not a copy. Writing back into it modifies the
        original frame in place, which is why the assignment on the second line is necessary rather
        than optional --- \code{addWeighted} returns a new array.
  \item The ROI is clamped to the image, and the function returns early if the clamp produces an
        empty region:
        \code{if x1 <= x0 or y1 <= y0: return frame}. Without that guard, a frame smaller than the
        card raises on the slice assignment. A test covers this case with a $20\times20$ frame.
\end{itemize}

\begin{edgetrap}{Fixed-size cards overflow with two hands}
The card is sized dynamically --- both dimensions --- from the text it must hold:

\begin{lstlisting}
sizes = [cv2.getTextSize(line, ...)[0] for line in lines]
text_width = max(width for width, _ in sizes)
line_height = text_height + 12
x1 = min(width, x0 + text_width + 2 * HUD_PAD)
y1 = min(height, y0 + 2 * HUD_PAD + line_height * len(lines))
\end{lstlisting}

A hard-coded box would work until the moment two hands are detected, at which point the HUD gains two
extra lines and the text spills out of the card --- over the live video, unreadable. Measuring the
text first makes the card correct for one hand, two hands, or a long ``Rock On'' gesture label.
\end{edgetrap}

\section{Teardown and resource safety}
\label{sec:teardown}

\begin{codecard}
\begin{lstlisting}
try:
    while True:
        ok, frame = capture.read()
        ...
finally:
    capture.release()
    tracker.close()
    cv2.destroyAllWindows()
\end{lstlisting}
\end{codecard}

The \code{finally} block runs on every path out of the loop: the user pressing \code{q}, an
\code{Esc}, a failed frame read that breaks the loop, or an exception. Each call releases a distinct
resource:

\begin{itemize}
  \item \code{capture.release()} returns the camera device to the operating system. Skipping it is
        what produces the ``camera already in use'' error on the next run.
  \item \code{tracker.close()} releases the MediaPipe graph and its native resources.
  \item \code{cv2.destroyAllWindows()} removes the GUI window. Without it, a crashed run can leave an
        orphaned window that keeps a stale frame on screen and is not obviously connected to any
        process.
\end{itemize}

\begin{keyidea}
\textbf{The general principle:} acquire resources in \code{\_\_init\_\_} or at the start of a
function, and release them in \code{finally} --- never at the end of the happy path only. The happy
path is the one case that is already working.
\end{keyidea}

\section{The full data flow}

\begin{figure}[h]
\centering
\begin{tikzpicture}[
  b/.style={draw=slateline, fill=slatebg, rounded corners=1mm, align=center,
            inner sep=1.6mm, font=\scriptsize\sffamily, text=navy,
            text width=2.4cm, minimum height=1.05cm},
  io/.style={b, draw=amberx, fill=ambersoft},
  pure/.style={b, draw=tealx, fill=tealsoft},
  ar/.style={-{Stealth[length=1.8mm]}, draw=cyanx, line width=0.8pt},
  back/.style={-{Stealth[length=1.8mm]}, draw=slatetext, line width=0.7pt, dashed}
]
  % ---- top row: left to right ----
  \node[io]   (cam)  at ( 1.30, 1.40) {camera frame\\BGR $480\times640\times3$};
  \node[b]    (flip) at ( 4.25, 1.40) {\code{cv2.flip}\\mirror};
  \node[b]    (rgb)  at ( 7.20, 1.40) {\code{cvtColor}\\BGR$\to$RGB};
  \node[b]    (palm) at (10.15, 1.40) {BlazePalm\\palm box};
  \node[b]    (lm)   at (13.10, 1.40) {landmark\\regression};

  % ---- bottom row: right to left ----
  \node[pure] (parse) at (13.10,-0.60) {\code{parse\_hand\_data}\\one-pass cache};
  \node[pure] (geom)  at (10.15,-0.60) {\code{fingers\_state}\\+ palm normal};
  \node[pure] (db)    at ( 7.20,-0.60) {\code{GestureDebouncer}};
  \node[b]    (hud)   at ( 4.25,-0.60) {alpha-blended\\HUD};
  \node[io]   (scr)   at ( 1.30,-0.60) {window\\(\code{imshow})};

  % ---- flow ----
  \draw[ar] (cam)  -- (flip);
  \draw[ar] (flip) -- (rgb);
  \draw[ar] (rgb)  -- (palm);
  \draw[ar] (palm) -- (lm);
  \draw[ar] (lm)   -- (parse);
  \draw[ar] (parse)-- (geom);
  \draw[ar] (geom) -- (db);
  \draw[ar] (db)   -- (hud);
  \draw[ar] (hud)  -- (scr);

  % ---- feedback loop, left-hand side ----
  \draw[back] (0.10,-0.60) -- (-0.80,-0.60) -- (-0.80,1.40) -- (0.10,1.40);
  \node[font=\scriptsize\sffamily, text=slatetext, rotate=90, anchor=south]
       at (-0.98, 0.40) {next frame, 30$\times$/s};

  % ---- legend ----
  \node[draw=amberx, fill=ambersoft, minimum size=3.4mm, inner sep=0pt] at (4.30,-1.90) {};
  \node[font=\scriptsize\sffamily, text=slatetext, anchor=west] at (4.62,-1.90)
       {I/O and the model (\code{hand\_tracker.py})};
  \node[draw=tealx, fill=tealsoft, minimum size=3.4mm, inner sep=0pt] at (4.30,-2.30) {};
  \node[font=\scriptsize\sffamily, text=slatetext, anchor=west] at (4.62,-2.30)
       {pure logic (\code{hand\_logic.py})};
\end{tikzpicture}
\caption{One frame, end to end. The trace runs left to right along the top row --- mirror, colour
conversion, detect, regress --- then drops into the pure-logic row and travels back right to left to
the window. The dashed line on the left is the only thing making this a video rather than a snapshot.}
\label{fig:pipeline}
\end{figure}

\begin{keyidea}
\textbf{Chapter summary.} The codebase is split along exactly one seam --- pure logic versus
everything that touches hardware --- and that seam is what makes the project testable at all. Each
frame converts its landmarks once into a per-hand dictionary, which every downstream rule reads.
Opening a camera searches both backends and indices before failing with a message that names the
likely cause, and all resources are released in a \code{finally} block. The \code{--image} mode runs
the identical pipeline on a still, which is what allows the whole program, drawing included, to be
verified in CI without a display server.
\end{keyidea}

\chapter{Complete Walkthrough of the Source Code}
\label{ch:source}

This chapter walks the three source files from top to bottom. Every code block below is
\textbf{quoted directly out of the project files with \LaTeX{}'s \code{lstinputlisting}}, so the line
numbers printed in the margin are the file's real line numbers --- open \code{hand_logic.py}, jump to
line 127, and you will be looking at exactly what is printed here. The complete files are reproduced
verbatim in Appendix~\ref{app:source}.

The three files are presented in dependency order: the pure logic first, then the application that
drives it, then the tests that pin both down.

\section{\code{hand\_logic.py} --- the pure layer}
\label{sec:src-logic}

\subsection{Module docstring, imports and constants}

\begin{codecard}
\lstinputlisting[firstline=1, lastline=30, firstnumber=1]{../hand_logic.py}
\end{codecard}

\begin{itemize}
  \item \textbf{Lines 1--9, the docstring.} It states two things a reader cannot infer from the code:
        the \emph{dependency policy} (``imports nothing but the standard library'') and the
        \emph{coordinate convention} (origin top-left, $y$ grows downward). The second is the single
        most consequential fact in the module --- every comparison below depends on it --- so it is
        declared before any code.
  \item \textbf{Line 11, \code{from \_\_future\_\_ import annotations}.} Makes type annotations lazy:
        they are stored as strings rather than evaluated. That matters because it lets a hint like
        \code{list[dict]} be written on Python 3.8 as well as 3.9+, without a runtime cost. It is the
        standard first line of any modern Python module.
  \item \textbf{Lines 13--14, the only two imports in the module.} \code{math} for \code{math.hypot},
        and \code{deque} for the debouncer's sliding window. No \code{cv2}, no \code{mediapipe} ---
        the dependency arrow from Chapter~\ref{ch:architecture} enforced by construction.
  \item \textbf{Lines 16--21, the landmark index constants.} These are \emph{named} indices. Writing
        \code{PINKY\_MCP} rather than \code{17} at the point of use is what allows
        Section~\ref{sec:palm} to be checked against Section~\ref{sec:src-logic} line by line: the
        code and the mathematics use the same names.
  \item \textbf{Line 23, \code{FINGER\_PAIRS}.} A list of (tip, PIP) tuples in the order the state
        list will use: index, middle, ring, pinky. Note that it deliberately \emph{excludes} the
        thumb, which follows a completely different rule.
  \item \textbf{Lines 25--30, the remaining constants.} \code{EXPECTED\_LANDMARKS = 21} is used as a
        validity check. \code{THUMB\_DISTANCE} and \code{THUMB\_X} are the string names of the two
        thumb rules --- string constants rather than booleans, so they read well in a CLI and in an
        error message. \code{NO\_GESTURE} is the placeholder emitted when a pose matches no pattern,
        and --- importantly --- the initial value of the debouncer, which is what defines its
        warm-up behaviour.
\end{itemize}

\subsection{The distance helper}

\begin{codecard}
\lstinputlisting[firstline=33, lastline=34, firstnumber=33]{../hand_logic.py}
\end{codecard}

\begin{itemize}
  \item \code{math.hypot(dx, dy)} computes $\sqrt{dx^2 + dy^2}$ without forming the squares
        explicitly, and without the overflow that naive squaring can cause for very large values.
        Here the coordinates are pixels, so overflow is not a concern --- but the substitution also
        avoids a temporary and reads exactly like the formula in Chapter~\ref{ch:math}.
  \item The leading underscore is a Python convention: \code{\_distance} is \emph{internal}. It is
        not imported by any test, signalling that callers outside the module should not depend on it.
\end{itemize}

\subsection{Normalized to pixel conversion, and the bounding box}

\begin{codecard}
\lstinputlisting[firstline=37, lastline=46, firstnumber=37]{../hand_logic.py}
\end{codecard}

\begin{itemize}
  \item \textbf{Line 39} is the comprehension from Chapter~\ref{ch:foundations}: for each normalized
        landmark, multiply $x$ by the frame width and $y$ by the frame height, and truncate to an
        integer with \code{int()}. Note the argument order \code{(x, y)} in the output tuple ---
        pixels are stored as $(x, y)$ pairs even though the NumPy array is indexed \code{[y, x]}.
        This module therefore speaks ``$(x, y)$'' consistently, and only the array indexing code in
        the application has to remember the difference.
  \item \textbf{Line 44--45} build two parallel lists of $x$ and $y$ coordinates. This is the
        standard idiom for reducing a point list to axis extents, and it is slightly clearer than a
        running min/max loop.
  \item \textbf{Line 46} returns the box in the conventional order
        $(x_{\min}, y_{\min}, x_{\max}, y_{\max})$. The order is fixed and documented, because a
        rectangle whose coordinates are swapped is a silent, plausible-looking bug.
\end{itemize}

\subsection{The palm normal}

\begin{codecard}
\lstinputlisting[firstline=49, lastline=61, firstnumber=49]{../hand_logic.py}
\end{codecard}

\begin{itemize}
  \item \textbf{Line 50--56, the docstring.} It writes the formula out in the comment as well as the
        code \emph{and} states the sign convention. For a function whose output is a bare float whose
        meaning is entirely in its sign, that is essential: without the docstring, the returned
        \code{-8775.0} is meaningless.
  \item \textbf{Lines 57--59} destructure the three landmarks used --- wrist, index MCP, pinky MCP
        --- into named local variables. This costs one line each and removes every possibility of
        confusing \code{x5} with \code{x17} while reading the formula.
  \item \textbf{Line 60} is the cross product's $z$-component exactly as derived in
        Section~\ref{sec:palm}. The parenthesisation is explicit rather than relying on operator
        precedence, because a subtle precedence error here would produce a plausible number with the
        wrong sign.
  \item \code{float(...)} on the return is defensive: heights and widths may arrive as NumPy
        integers, and \code{np.int64 * np.int64} can overflow at 32 bits. Converting to Python's
        arbitrary-precision float sidesteps that entirely.
\end{itemize}

\begin{codecard}
\lstinputlisting[firstline=64, lastline=66, firstnumber=64]{../hand_logic.py}
\end{codecard}

A three-line function that exists to give a sign a name. It is separate from \code{palm\_normal\_z}
so that the \emph{convention} lives in exactly one place: if the camera were ever un-mirrored, this
single line is what would change.

\subsection{Reading the handedness label safely}

\begin{codecard}
\lstinputlisting[firstline=69, lastline=78, firstnumber=69]{../hand_logic.py}
\end{codecard}

\begin{itemize}
  \item \textbf{Lines 70--73}: the attribute chain \code{handedness.classification[0].label} can fail
        in three distinct ways --- the object may lack the attribute, the list may be empty, or the
        value may not be subscriptable at all. Catching all three exception types in one clause is
        correct here because the recovery is identical for each: the label is unknown.
  \item \textbf{Lines 74--75}: an unexpected string is normalised to \code{"Unknown"} rather than
        passed through. Everything downstream then has exactly three possible labels, which is what
        makes the \code{in ("Left", "Right")} checks elsewhere total.
  \item \textbf{Lines 76--77}: \code{--swap-handedness} is applied \emph{here}, at parse time, so
        that the swapped label is what the rest of the program sees --- including the state keys used
        for smoothing and debouncing. Swapping later (at draw time) would have left the state keyed
        by the original, inconsistent label.
\end{itemize}

\subsection{Parsing a frame: the single-pass cache}

\begin{codecard}
\lstinputlisting[firstline=81, lastline=124, firstnumber=81]{../hand_logic.py}
\end{codecard}

This is the most important function in the module, so it is worth reading line by line.

\begin{itemize}
  \item \textbf{Lines 82--92, the docstring.} It states the duck-typing contract explicitly
        (``only ever touches \code{multi\_hand\_landmarks} and \code{multi\_handedness}''), which is
        the reason the tests can pass a fake object and get real coverage. It also documents the
        exact dictionary schema --- six keys --- and the two invariants: never raises, and skips
        incomplete hands.
  \item \textbf{Lines 93--95}, \code{getattr(results, "multi\_hand\_landmarks", None)} and the
        truthiness check. A result with no hands may or may not have the attribute; either way,
        \code{None} or an empty list returns \code{[]}. This one \code{if} is what makes the
        ``zero hands'' and ``\code{None} input'' tests pass without a special case.
  \item \textbf{Lines 97--99}: \code{or []} converts a present-but-\code{None} attribute into an
        empty list, so the length check on the next line is always safe.
  \item \textbf{Lines 101--104, the loop}. \code{enumerate} supplies both the index (used to look up
        the matching handedness entry) and the hand itself. The index correspondence matters:
        MediaPipe emits the two lists in the same order, which is the only reason
        \code{handedness[index]} is correct.
  \item \textbf{Lines 105--107, the completeness gate.} A hand without a full set of 21 landmarks is
        skipped with \code{continue} rather than \code{break}: the \emph{other} hand in the frame may
        be perfectly usable, and discarding it would be a worse outcome than ignoring the malformed
        one.
  \item \textbf{Lines 109--110}: \code{norm} keeps the $(x, y, z)$ triples, and \code{px} converts
        with the same helper described above. Both are stored because they serve different consumers
        --- \code{norm} for debugging and future 3D work, \code{px} for every rule in the module.
  \item \textbf{Lines 111--124, the dictionary literal.} Every derived value is computed \emph{here
        and once}: the palm normal, the orientation label, and the bounding box. This is the
        single-pass cache of Chapter~\ref{ch:architecture} --- the reason the render loop performs no
        redundant coordinate arithmetic.
\end{itemize}

\subsection{The finger rules}

\begin{codecard}
\lstinputlisting[firstline=127, lastline=162, firstnumber=127]{../hand_logic.py}
\end{codecard}

\begin{itemize}
  \item \textbf{Lines 128--139, the docstring.} It documents all three rules --- the vertical test for
        the four long fingers, the distance rule for the thumb, and the alternative horizontal rule
        --- plus the exact safe-result contract. It also explains why \code{palm\_z} is accepted but
        unused by the default rule (API symmetry with the orientation readout), which prevents a
        future reader from assuming the parameter is a bug.
  \item \textbf{Lines 141--142, the input guard.} This single line implements the specification's
        requirement: a \code{None}, empty, or short landmark list returns
        \code{[False] * 5} and never raises. Note it is a \emph{list comprehension-free} early
        return, evaluated before any indexing can occur.
  \item \textbf{Line 146, the four-finger test.} A single comprehension over \code{FINGER\_PAIRS},
        comparing the $y$ of each tip against the $y$ of its PIP joint --- the rule derived in
        Section~\ref{sec:screencs}. The result is a list of exactly four booleans, in the order the
        pairs are declared.
  \item \textbf{Lines 148--150, the horizontal thumb rule.} Two conditions guard it: an explicit
        \code{thumb\_mode == THUMB\_X} \emph{and} a usable label. If either fails, control falls to
        the \code{else} branch --- so an \code{"Unknown"} label degrades to the distance rule instead
        of silently reporting a folded thumb. That fallback is asserted by a test.
  \item \textbf{Lines 151--154, the distance rule.} The default path, and the return of
        Section~\ref{sec:thumb}'s derivation. Written as a direct transcription of the inequality,
        which is why it is readable without a comment.
  \item \textbf{Lines 155--158, the exception guard.} The commented contract is ``never raises'',
        and a comprehension over user-supplied data cannot guarantee that. Catching
        \code{IndexError}, \code{TypeError} and \code{ValueError} covers short lists, \code{None}
        elements and non-numeric entries, and all three recover to the same safe value. A test feeds
        a list of 21 \code{None} values specifically to exercise this path.
  \item \textbf{Line 160}: the thumb result is \emph{prepended} to the four finger results, giving
        the canonical \code{[thumb, index, middle, ring, pinky]} order that every consumer ---
        classification, HUD text, and the tests --- relies on.
\end{itemize}

\subsection{Gesture classification}

\begin{codecard}
\lstinputlisting[firstline=165, lastline=182, firstnumber=165]{../hand_logic.py}
\end{codecard}

\begin{itemize}
  \item \textbf{Lines 166--167}: a short or empty state list returns the placeholder rather than
        raising, keeping the module's no-exceptions contract intact.
  \item \textbf{Line 169} unpacks exactly five values and coerces each with \code{bool()}. The
        coercion matters: it means a caller passing \code{1}/\code{0}, or NumPy booleans, gets the
        same behaviour as one passing \code{True}/\code{False}.
  \item \textbf{Lines 171--180, the pattern chain.} Each test is \emph{exhaustive} --- it states
        which fingers must be down as well as which must be up. Without the negated terms, ``Peace''
        would also match an open palm, and the first match would win. The ordering from most to least
        obvious is a readability choice, not a correctness requirement, because the patterns are
        mutually exclusive.
  \item \textbf{Line 182}: anything unmatched returns \code{NO\_GESTURE}. The module never returns
        \code{None} for a gesture, so the HUD can always print something.
\end{itemize}

\subsection{The debouncer}

\begin{codecard}
\lstinputlisting[firstline=185, lastline=207, firstnumber=185]{../hand_logic.py}
\end{codecard}

\begin{itemize}
  \item \textbf{Lines 185--191, the class docstring.} It states the contract in three clauses: the
        window length, the agreement threshold, and --- crucially --- the warm-up behaviour
        (``before any majority has ever formed it returns \code{---}''). That last sentence is the
        answer to a question the code alone leaves ambiguous.
  \item \textbf{Lines 193--198, the constructor.} Three pieces of state per instance:
        \code{window} and \code{min\_agree} as configuration, \code{\_history} as a
        \code{deque(maxlen=window)} that evicts automatically, and \code{\_stable} seeded with the
        placeholder. Seeding with the placeholder \emph{is} the warm-up policy.
  \item \textbf{Lines 200--201, the \code{stable} property.} A read-only view of the last accepted
        label. Exposing it via \code{@property} rather than a plain attribute means callers cannot
        accidentally assign to it --- a small invariant guard.
  \item \textbf{Lines 203--207, \code{update}.} Three statements: append the new observation (the
        \code{deque} silently evicts the oldest), test whether the \emph{new} label alone has reached
        the threshold within the window, and return the current stable label either way. Counting
        only the newest label is correct, and cheaper than building a histogram of all five entries:
        if any label reached the threshold, the newest one is the only candidate that could have just
        crossed it.
\end{itemize}

\subsection{The exponential moving average}

\begin{codecard}
\lstinputlisting[firstline=210, lastline=220, firstnumber=210]{../hand_logic.py}
\end{codecard}

\begin{itemize}
  \item \textbf{Lines 211--214, the docstring}, which writes the recurrence in words and states the
        first-sample rule.
  \item \textbf{Lines 215--216}: the \code{None} special case. Without it, the arithmetic would
        raise, or --- worse, if a caller substituted zero --- the dot would animate in from the
        origin on the first frame.
  \item \textbf{Lines 217--220}: the weighted sum, applied to $x$ and $y$ independently but in one
        expression, returning a plain tuple. Floats are returned, not integers: rounding here would
        accumulate error and defeat the purpose of smoothing. The caller rounds once, at draw time.
\end{itemize}

\section{\code{hand\_tracker.py} --- the application}
\label{sec:src-tracker}

\subsection{Module setup, the import guard, and constants}

\begin{codecard}
\lstinputlisting[firstline=36, lastline=71, firstnumber=36]{../hand_tracker.py}
\end{codecard}

\begin{itemize}
  \item \textbf{Line 36}: the module logger. Every diagnostic in the application goes through it
        rather than \code{print}, so verbosity is controllable from the command line and library
        users are not spammed by a chatty dependency.
  \item \textbf{The \code{warnings.catch\_warnings()} block} around the MediaPipe import suppresses
        the deprecation noise of the legacy API --- scoped to the import, so real warnings elsewhere
        still surface. This is discussed in detail in Section~\ref{sec:deps}.
  \item \textbf{The \code{RuntimeError} guard} is the third dependency trap from
        Chapter~\ref{ch:engine}: it converts an obscure \code{AttributeError} from deep inside a
        library into one actionable sentence naming the fix.
  \item \textbf{The \code{tikz}-free constants block (lines 54--71)} collects every magic number the
        application uses: the window title, the captures directory, the default frame size, the six
        HUD styling values, the two index-tip marker radii and its colour, and the FPS averaging
        window. Collecting them means the look of the overlay can be tuned without hunting through
        the drawing code --- and it makes each one citable in this handbook.
  \item \code{INDEX\_TIP\_COLOR = (0, 165, 255)} deserves a note: this is \textbf{BGR}, not RGB.
        Read as a colour, it is ``no blue, 165 green, 255 red'' --- the orange dot visible in the
        running application. A constant written as \code{(255, 165, 0)} would draw a \emph{blue}
        marker, and it would look deliberate.
\end{itemize}

\subsection{\code{HandTracker}: construction and the detection call}

\begin{codecard}
\lstinputlisting[firstline=77, lastline=122, firstnumber=77]{../hand_tracker.py}
\end{codecard}

\begin{itemize}
  \item \textbf{Lines 77--96, the constructor.} The five parameters mirror the CLI flags, each with
        the default the specification requires. Four of them are configuration; the fifth,
        \code{mode}, is passed straight to \code{static\_image\_mode}. Note the ordering of
        assignments: the three caches are created before the MediaPipe model, so an object is never
        in a half-built state even if model construction raises.
  \item \textbf{Line 90, the model.} This is the only place in the program where MediaPipe is
        instantiated. Isolating construction to one line makes the three thresholds from
        Chapter~\ref{ch:engine} easy to find and tune.
  \item \textbf{Line 91--93, the three caches.} \code{hands\_data} is the per-frame telemetry;
        \code{\_smooth\_tip} holds one EMA state per hand key; \code{\_debounce} holds one
        \code{GestureDebouncer} per hand key. All three are \emph{instance} attributes --- two
        trackers would not share state.
  \item \textbf{Lines 98--114, the detection itself.} The conversion to RGB, the read-only flag, the
        model call, and the restore of the write flag. Setting \code{rgb.flags.writeable = False} is
        a micro-optimisation that tells NumPy the buffer will not be written, letting the model read
        it without defensive copying.
  \item \textbf{Lines 116--119, the cache rebuild.} Note the order: the image is measured
        (\code{frame.shape[:2]} gives \emph{height, width}), and \code{parse\_hand\_data} is called
        with \code{(width, height)} --- the opposite order, because the function's parameters are
        named \code{width, height}. This is the single most error-prone line in the file, and it is
        correct precisely because both names are explicit at the call site.
  \item \textbf{Lines 121--122, the optional mesh.} Drawing is guarded by both the \code{draw} flag
        and the presence of landmarks. \code{mp\_drawing.draw\_landmarks} is handed two styles: a
        green joint style and a near-white connection style, which is what produces the familiar
        MediaPipe skeleton.
\end{itemize}

\subsection{The accessor methods}

\begin{codecard}
\lstinputlisting[firstline=124, lastline=147, firstnumber=124]{../hand_tracker.py}
\end{codecard}

\begin{itemize}
  \item \textbf{\code{hand\_count}} is a one-liner over the cache. It is a method rather than a
        direct attribute read so that the length and the cache can never disagree.
  \item \textbf{\code{hand\_label}} returns \code{"Unknown"} out of range --- never raises, and never
        returns \code{None}, so its result can be placed directly into HUD text.
  \item \textbf{\code{find\_positions}} is the specification's method, preserved exactly. Two things
        about it are deliberate: it returns a \emph{copy} (\code{list(...)}) so a caller cannot
        mutate the cache, and it returns \code{[]} --- not \code{None} --- for a bad index, so the
        finger rules receive a value they already have a guard for. The docstring records the
        vestigial \code{frame} parameter and the size caveat from Section~\ref{sec:cache}.
  \item \textbf{\code{fingers\_up}} is a one-line delegation to \code{fingers\_state}, carrying the
        instance's \code{thumb\_mode} so the CLI flag reaches the rule without the render loop having
        to pass it around.
\end{itemize}

\subsection{The smoothed tip, orientation, and debounced gesture}

\begin{codecard}
\lstinputlisting[firstline=149, lastline=178, firstnumber=149]{../hand_tracker.py}
\end{codecard}

\begin{itemize}
  \item \textbf{Lines 150--152, \code{index\_tip}'s guards.} Out-of-range index and a too-short
        landmark list both return \code{None}. \code{None} is the right sentinel here --- unlike
        \code{find\_positions}, the caller genuinely has nothing to draw --- and the render loop
        checks for it before calling \code{cv2.circle}.
  \item \textbf{Line 157, the state key.} The EMA state is looked up by hand identity, not by list
        position, which is the fix for the cross-contamination trap in Section~\ref{sec:ema}.
  \item \textbf{Lines 158--160}: the smoothing call, the state update, and the cast to integers. Note
        that the \emph{float} value is what gets stored, and the rounding happens only on the way out
        --- so the filter's memory is never quantised.
  \item \textbf{\code{orientation}} returns the cached label, with an em-dash for out-of-range. It
        performs no arithmetic at all: the sign test was done once, in \code{parse\_hand\_data}.
  \item \textbf{Lines 169--178, \code{gesture}.} The debouncer is created lazily on first use for a
        given key and then reused. Lazy creation is what allows the tracker to adapt to hands
        appearing and disappearing mid-session without a separate bookkeeping pass --- and because
        the key is the label, a hand that briefly leaves the frame and returns resumes its own
        gesture history rather than inheriting another hand's.
\end{itemize}

\begin{codecard}
\lstinputlisting[firstline=179, lastline=189, firstnumber=179]{../hand_tracker.py}
\end{codecard}

\begin{itemize}
  \item \textbf{\code{\_state\_key}} is a static method because it needs no instance state --- it is
        a pure function of its arguments, and marking it \code{@staticmethod} says so.
  \item The fallback key, \code{f"index:\{hand\_index\}"}, is deliberately namespaced with a prefix.
        If the label list ever contained the literal string \code{"0"}, the two key spaces could
        collide; the prefix makes that impossible.
  \item \textbf{\code{close}} releases the MediaPipe graph. It exists so the teardown in
        Section~\ref{sec:teardown} can call one method rather than reaching into
        \code{self.hands}, keeping the model's lifecycle inside the class that owns it.
\end{itemize}

\subsection{The HUD: text assembly and drawing}

\begin{codecard}
\lstinputlisting[firstline=192, lastline=230, firstnumber=192]{../hand_tracker.py}
\end{codecard}

\begin{itemize}
  \item \textbf{Lines 192--194, \code{hud\_lines}.} A pure function that assembles the text: FPS,
        total, then one line per hand. Separating assembly from drawing means the text can be tested
        without a frame --- and it is, by \code{test\_hud\_lines\_formatting}.
  \item \textbf{Lines 198--199}: an empty line list returns immediately, so an empty card is never
        drawn.
  \item \textbf{Lines 201--205, measurement.} \code{cv2.getTextSize} returns the pixel extent of a
        string in the chosen font. Measuring \emph{before} drawing is what allows the card to size
        itself to its content, which is the fix for the two-hand overflow trap in
        Section~\ref{sec:hud}.
  \item \textbf{Lines 207--214, the clamped rectangle.} Four values are computed: the top-left
        corner, and the bottom-right corner clamped to the image with \code{min}. The early return on
        \code{x1 <= x0 or y1 <= y0} handles a frame so small that the card would be degenerate ---
        20-pixel-square frames do occur in tests and in thumbnail previews.
  \item \textbf{Lines 216--217, the blend.} The ROI is a NumPy \emph{view}; \code{addWeighted} reads
        it, and the assignment writes the darkened result back into the same memory. The expression
        \code{1.0 - HUD\_ALPHA} keeps the two weights summing to 1, so brightness is scaled rather
        than shifted --- which is why the card looks like tinted glass rather than a black box.
  \item \textbf{Lines 219--230, the text loop.} The baseline advances by \code{line\_height} each
        iteration, starting one pad plus one text height below the card's top edge. White
        (\code{255, 255, 255}) is chosen because the blended background is guaranteed dark; the
        \code{cv2.LINE\_AA} flag enables anti-aliasing so small text stays legible.
\end{itemize}

\subsection{The per-frame loop body}

\begin{codecard}
\lstinputlisting[firstline=233, lastline=254, firstnumber=233]{../hand_tracker.py}
\end{codecard}

This function is the shared heart of both the live loop and \code{--image} mode --- which is what
guarantees the still-image path exercises the real drawing code rather than a copy of it.

\begin{itemize}
  \item \textbf{Line 238, the loop bound.} \code{range(tracker.hand\_count())} makes both hands in
        the frame fall out naturally, and makes the two-hand requirement of the specification a
        non-event.
  \item \textbf{Lines 239--242, the read-and-read-again pattern.} \code{hand\_label} and
        \code{find\_positions} both come from the cache populated by the preceding
        \code{find\_hands}. No landmark is recomputed below this line.
  \item \textbf{Lines 243--244}: \code{sum(states)} counts the \code{True} values directly. Since the
        booleans are Python \code{bool}, which is a subclass of \code{int}, this is both correct and
        idiomatic.
  \item \textbf{Line 245, the HUD row.} An f-string with four fields: label, count, debounced
        gesture, and the palm orientation in square brackets. This single line is the ``diagnostic
        visual telemetry'' the specification asks for, and every field is traceable to a rule
        derived in Chapter~\ref{ch:math}.
  \item \textbf{Lines 247--251, the marker.} Two filled circles: a white halo of radius 14, then the
        orange dot of radius 10 on top. Drawing the halo first and the dot second is what produces a
        ringed marker that stays visible against both light and dark backgrounds --- the same design
        principle as the HUD card.
\end{itemize}

\subsection{Screenshots, and the camera fallback chain}

\begin{codecard}
\lstinputlisting[firstline=257, lastline=289, firstnumber=257]{../hand_tracker.py}
\end{codecard}

\begin{itemize}
  \item \textbf{\code{save\_screenshot}} (lines 257--265) creates the directory if needed --- the
        \code{exist\_ok=True} flag makes repeated presses idempotent rather than errors --- and builds
        a timestamped filename. The timestamp uses \code{\%f} microseconds truncated to milliseconds
        (\code{[:-3]}), so two presses within the same second cannot overwrite each other.
  \item The \code{if cv2.imwrite(...)} test is not paranoia: \code{imwrite} returns \code{False} on
        failure (a full disk, a permission problem, an unsupported extension) rather than raising, so
        an unchecked call would fail silently and the user would simply never find the file.
  \item \textbf{\code{open\_camera}} (lines 268--289) is the fallback chain of
        Section~\ref{sec:camera}, covered in detail there. Two details are worth re-reading in the
        code: the backend list is chosen by \code{os.name == "nt"}, and the candidate list is built
        by filtering the requested index out of \code{(1, 0)} so no index is ever tried twice.
\end{itemize}

\subsection{The live capture loop}

\begin{codecard}
\lstinputlisting[firstline=292, lastline=338, firstnumber=292]{../hand_tracker.py}
\end{codecard}

\begin{itemize}
  \item \textbf{Lines 294--296}: the camera is opened \emph{first}, and a \code{None} return exits
        with status 1 before any model is constructed. Failing fast here means a missing webcam never
        pays the several-hundred-millisecond cost of loading the MediaPipe graph.
  \item \textbf{Lines 304--305}: a \code{deque(maxlen=30)} of instantaneous frame rates and a
        timestamp. The window gives a one-second moving average at 30 FPS --- long enough to smooth
        the reading, short enough to react to a real slowdown.
  \item \textbf{Lines 314--316, the per-frame core.} Flip, detect, annotate. Three lines, and
        \code{annotate\_hands} is shared with \code{--image} mode.
  \item \textbf{Lines 318--323, the FPS calculation.} This is the \code{time.time()} delta approach
        the specification requires: measure the interval since the previous frame, invert it for an
        instantaneous rate, and average the last thirty. The \code{if delta > 0} guard protects
        against a clock with insufficient resolution returning the same timestamp twice --- which
        would otherwise divide by zero.
  \item \textbf{Lines 325--326}: the HUD is drawn and the frame is shown. The order matters ---
        \code{imshow} displays whatever is in the buffer at that moment, so drawing after it would
        show a frame-delayed overlay.
  \item \textbf{Lines 328--332, input handling.} \code{cv2.waitKey(1)} waits one millisecond for a
        keypress, which also pumps the GUI event loop --- without it the window would freeze.
        \code{\& 0xFF} masks off platform-specific bits above the ASCII range, so the comparison
        against \code{ord("q")} works on every platform. Key 27 is \code{Esc}.
  \item \textbf{Lines 333--336, the \code{finally} block}: camera released, model closed, windows
        destroyed --- on every exit path, as required by Section~\ref{sec:teardown}.
\end{itemize}

\subsection{Headless image mode}

\begin{codecard}
\lstinputlisting[firstline=341, lastline=373, firstnumber=341]{../hand_tracker.py}
\end{codecard}

\begin{itemize}
  \item \textbf{Lines 342--346}: \code{cv2.imread} returns \code{None} for a missing or unreadable
        file --- it does not raise. Checking for \code{None} and returning 1 is what makes
        \code{--image does-not-exist.png} a clean failure that a CI script can detect from the exit
        code.
  \item \textbf{Lines 348--350}: resizing happens only when \emph{both} dimensions were supplied.
        The default of ``native resolution'' is the right choice for stills, where there is no
        frame-rate argument for downscaling.
  \item \textbf{Lines 352--357}: the tracker is constructed with \code{mode=True}
        (\code{static\_image\_mode}), because a single frame has no successor to track into.
  \item \textbf{Lines 359--363, the shared pipeline.} \code{find\_hands} $\to$
        \code{annotate\_hands} $\to$ \code{draw\_hud} --- the identical functions the live loop uses.
        The \code{try/finally} closes the model even if detection raises.
  \item \textbf{Lines 365--373, the output path.} \code{os.path.splitext} splits the input into stem
        and extension; the stem plus \code{\_out.png} is written \emph{next to the input}, so running
        the mode twice is idempotent and no directory configuration is needed. The reported
        \code{hand\_count()} comes from the cache built during \code{find\_hands}, which is why the
        log line can state a hand count without re-running detection.
\end{itemize}

\subsection{The command-line interface and entry point}

\begin{codecard}
\lstinputlisting[firstline=376, lastline=422, firstnumber=376]{../hand_tracker.py}
\end{codecard}

\begin{itemize}
  \item \textbf{Lines 377--381}: \code{argparse} is configured with a program name and a description,
        so \code{--help} produces usable output rather than a bare usage string.
  \item \textbf{Lines 382--384, the camera and size flags.} \code{--width} and \code{--height}
        default to \code{None}, not to 640 and 480. That deliberate choice is what allows
        \code{--image} mode to distinguish ``the user asked for a resize'' from ``the user accepted a
        default''; the live loop supplies 640$\times$480 itself when the flags are absent.
  \item \textbf{Line 388, \code{--thumb-mode}} uses \code{choices=("distance", "x")}. argparse then
        rejects a typo such as \code{--thumb-mode distancey} with a clear message, rather than
        silently falling back to the default.
  \item \textbf{Line 403, \code{--image}}: the flag that makes the whole application testable in CI,
        described in Section~\ref{sec:headless}.
  \item \textbf{Line 412--419, \code{main}}. Logging is configured from the parsed level, and the
        function returns an integer --- 0 or 1 --- rather than calling \code{sys.exit} itself. That
        makes \code{main()} callable from a test with an argument list, which is how a CLI is
        normally tested.
  \item \textbf{Lines 421--422, the entry point.} \code{raise SystemExit(main())} converts the return
        value into the process exit code, and the \code{\_\_name\_\_} guard ensures that importing
        the module (as the test suite does) never starts a camera.
\end{itemize}

\section{\code{tests/test\_hand\_logic.py} --- the offline test suite}
\label{sec:src-tests}

\subsection{The synthetic hands}

The file's first seventy lines define two coordinate sets by hand: \code{OPEN\_HAND} (all five
fingers extended) and \code{FIST} (all five folded). They are ordinary lists of $(x, y)$ tuples with
a comment naming each landmark. Because they are literals rather than captured data, every assertion
below is exact --- there is no tolerance and no randomness.

Four helpers manipulate them:

\begin{codecard}
\lstinputlisting[firstline=80, lastline=113, firstnumber=80]{../tests/test_hand_logic.py}
\end{codecard}

\begin{itemize}
  \item \code{\_with\_finger} and \code{\_with\_thumb} (lines 80--89) return a modified \emph{copy}
        of a hand with one finger, or the whole thumb, moved to a chosen position. Copying is
        essential: mutating the shared literal would make the tests order-dependent, so a test that
        passed alone could fail when run after another.
  \item \code{\_rotate} (lines 92--104) applies a real rotation matrix about the point set's
        centroid, rounding the result back to integers. It is the machinery behind the project's most
        important regression test --- the one that encodes the reviewer's original complaint about
        the horizontal thumb rule.
  \item \code{\_mirror\_x} (lines 107--113) negates every $x$ coordinate. Geometrically that is
        ``show the other side of the hand'', and it is exactly the transformation that must flip the
        palm normal's sign.
\end{itemize}

\subsection{The rotation-invariance test --- the project's key regression}

\begin{codecard}
\lstinputlisting[firstline=173, lastline=177, firstnumber=173]{../tests/test_hand_logic.py}
\end{codecard}

This four-line test is the reason the default thumb rule exists. It rotates the same open hand
through four angles and asserts that the thumb state never changes. The old horizontal rule fails
this test at 90\textdegree\ and 180\textdegree; the distance rule passes at every angle, for the
algebraic reason proved in Section~\ref{sec:thumb}. Turning a reviewer's complaint into an executable
assertion is what stops the bug from ever returning silently.

\subsection{The safe-result contract}

\begin{codecard}
\lstinputlisting[firstline=186, lastline=197, firstnumber=186]{../tests/test_hand_logic.py}
\end{codecard}

Four malformed inputs --- \code{None}, an empty list, a list one element too short, and a list of 21
\code{None} values --- all asserted to return exactly \code{[False] * 5}. The final line re-tests a
valid hand, guarding against the possibility that a guard was written so aggressively that it
suppresses correct input too. That is the shape of a complete boundary test: every invalid case
\emph{and} one valid case in the same test, so a regression in either direction fails.

\subsection{The fake MediaPipe results}

\begin{codecard}
\lstinputlisting[firstline=298, lastline=330, firstnumber=298]{../tests/test_hand_logic.py}
\end{codecard}

These five tiny classes are the reason \code{parse\_hand\_data} is testable without MediaPipe
installed. Each mimics exactly one level of the real result object's attribute chain:
\code{FakeResults} holds the two lists, \code{FakeHand} wraps a landmark list,
\code{FakeClassificationList} wraps a classification, and \code{FakeClassification} carries the
label. The full dotted chain --- from \code{multi\_handedness} through \code{classification[0]} down
to \code{label} --- therefore resolves exactly as it does in production, and the duck-typed parser
cannot tell the difference.

\code{FakeResults} has one particularly useful parameter:

\begin{itemize}
  \item \code{include\_handedness=False} omits the attribute \emph{entirely} on construction, which
        is how the ``missing \code{multi\_handedness}'' branch is covered. Passing an empty list
        would test something different: an attribute that exists but is empty. Both are real
        scenarios, and both are tested.
\end{itemize}

\subsection{Parse tests}

\begin{codecard}
\lstinputlisting[firstline=333, lastline=363, firstnumber=333]{../tests/test_hand_logic.py}
\end{codecard}

\begin{itemize}
  \item \textbf{\code{test\_parse\_two\_hands}} asserts on the label list, the landmark counts, the
        three-element \code{norm} tuples, the sign of \code{palm\_z}, and the orientation string ---
        then compares every pixel coordinate against the source with a $\pm 1$ tolerance. The
        tolerance is not sloppiness: the fake landmarks are built by dividing by 1000 and the parser
        multiplies by 1000 and truncates, so a half-unit difference is a floating-point artefact of
        the round trip, not a logic error. A genuinely wrong scale factor would be off by far more
        than one pixel.
  \item \textbf{\code{test\_parse\_zero\_hands}} covers the empty-but-present lists.
  \item \textbf{\code{test\_parse\_none\_and\_empty\_objects}} covers three distinct degenerate
        inputs: \code{None}, a bare \code{object()} with no attributes at all, and a result with
        neither list. All three must return \code{[]} rather than raising --- the ``never raises''
        contract is what allows the render loop to be written without a defensive \code{try} around
        it.
\end{itemize}

\begin{codecard}
\lstinputlisting[firstline=387, lastline=400, firstnumber=387]{../tests/test_hand_logic.py}
\end{codecard}

\begin{itemize}
  \item \textbf{\code{test\_parse\_skips\_incomplete\_hand}} builds a two-hand result where the first
        hand has only ten landmarks, and asserts that exactly one hand survives --- with the
        \emph{second} hand's label, proving the skip did not shift the handedness correspondence.
  \item \textbf{\code{test\_parse\_returns\_mirrored\_and\_back\_facing\_cases}} is the end-to-end
        confirmation of Section~\ref{sec:palm}: the same hand, mirrored in $x$, must flip the
        orientation label. It ties the geometry, the parser and the sign convention together in one
        assertion.
\end{itemize}

\begin{keyidea}
\textbf{Chapter summary.} The pure module converts coordinates once, applies three geometric rules
and two temporal filters, and never touches hardware or raises on bad input. The application layer
owns the model, the camera fallback chain, the HUD and the CLI, and routes both its live and headless
paths through the same annotation and drawing functions. The tests encode the two hardest-won
lessons of the project --- rotation invariance for the thumb, and a total, exception-free contract
for malformed landmarks --- as executable assertions rather than comments.
\end{keyidea}

\appendix

\chapter{Complete Source Listings}
\label{app:source}

The three files that make up the project, printed in full. These are the exact files from the
repository --- included here with \code{\textbackslash lstinputlisting}, not retyped --- so the line
numbers match those used throughout Chapter~\ref{ch:source}.

\section{\code{hand\_logic.py}}

The pure-logic module: no OpenCV, no MediaPipe, no camera. Everything in
Chapter~\ref{ch:math} lives here.

\lstinputlisting{../hand_logic.py}

\clearpage
\section{\code{hand\_tracker.py}}

The application: the model, the capture loop, the HUD, the CLI, the camera fallback chain, and the
headless \code{--image} mode.

\lstinputlisting{../hand_tracker.py}

\clearpage
\section{\code{tests/test\_hand\_logic.py}}

The offline suite. Synthetic coordinates only --- it runs in a bare Python interpreter with no
camera and no third-party packages installed.

\lstinputlisting{../tests/test_hand_logic.py}

\chapter{Verification Checklist and Troubleshooting}
\label{app:verify}

\section{Setup and first run}

\begin{center}
\rowcolors{2}{slatebg}{white}
\begin{tabularx}{\textwidth}{@{}r l X@{}}
\toprule
\rowcolor{white}
\textbf{\#} & \textbf{Step} & \textbf{What success looks like} \\
\midrule
1 & \code{python -m venv .venv} & A \code{.venv} directory appears. \\
2 & activate, then \code{pip install -r requirements.txt} & Three packages install:
  \code{mediapipe 0.10.14}, \code{opencv-contrib-python 4.11.0.86}, \code{numpy 1.26.x}. \\
3 & \code{pytest -q} & All tests pass. No camera is touched. \\
4 & \code{python hand\_tracker.py --image tests/fixtures/hand.png} & Prints a hand count and writes
  \code{tests/fixtures/hand\_out.png}; exit code 0. \\
5 & \code{python hand\_tracker.py} & A window titled \emph{Hand Tracker} opens with the mirrored
  feed. \\
\bottomrule
\end{tabularx}
\end{center}

\section{Live-mode acceptance tests}

\begin{center}
\rowcolors{2}{slatebg}{white}
\begin{tabularx}{\textwidth}{@{}l X l@{}}
\toprule
\rowcolor{white}
\textbf{Test} & \textbf{Procedure} & \textbf{Expected} \\
\midrule
Single hand & Raise one hand, fingers up. & HUD shows one line; the count matches your fingers. \\
Two hands & Raise both hands. & Two lines, and the total is the sum. \\
\textbf{Wrist rotation} & Rotate your wrist and turn the back of your hand toward the camera.
  & \textbf{The thumb count must not flip.} This is the fix the distance rule exists for. \\
Marker stability & Hold still and watch the index-finger dot. & It sits steady; no visible
  vibration. \\
Gesture debounce & Make a Peace sign, then a Fist. & The label changes once, after a few frames ---
  no flicker through intermediate states. \\
Empty scene & Remove your hands; cover the lens. & Count goes to 0; FPS keeps updating; no crash. \\
Screenshot & Press \code{s}. & A timestamped PNG appears in \code{captures/} with the HUD burned in. \\
Clean exit & Press \code{q}. & Window closes, camera light goes off, prompt returns. \\
Fallback & \code{python hand\_tracker.py --camera 5} & Logs a warning per failed index, then either
  opens another camera or exits with a clear message. \\
\bottomrule
\end{tabularx}
\end{center}

\section{Troubleshooting}

\begin{center}
\rowcolors{2}{slatebg}{white}
\begin{tabularx}{\textwidth}{@{}l X@{}}
\toprule
\rowcolor{white}
\textbf{Symptom} & \textbf{Cause and fix} \\
\midrule
\code{Could not open a webcam} & Another application holds the device (video call, browser tab, or a
  previous run that was killed). Close it and retry, or use \code{--camera 1}. \\
\code{RuntimeError: mediapipe>=0.10.30 removed mp.solutions} & MediaPipe was upgraded. Reinstall with
  \code{pip install -r requirements.txt}. \\
\code{ImportError: libGL.so.1} (Linux) & Missing system libraries for the non-headless OpenCV wheel:
  \code{sudo apt-get install -y libgl1 libglib2.0-0}. \\
Handedness labels look inverted & Expected on a mirrored feed. Display-only: relaunch with
  \code{--swap-handedness}. \\
Thumb never detects with \code{--thumb-mode x} & That rule assumes an upright, palm-facing hand. Use
  the default \code{distance} mode. \\
The dot is drawn in the wrong place & The frame changed size between \code{find\_hands} and
  \code{find\_positions}; the cache is valid for one frame size only
  (Section~\ref{sec:cache}). \\
Low FPS & Reduce \code{--width}/\code{--height}, or pass \code{--no-draw} to skip the mesh. \\
\bottomrule
\end{tabularx}
\end{center}

\chapter{Glossary and Quick Reference}
\label{app:glossary}

\section{Glossary}

\begin{description}
  \item[Alpha blending] Combining two images by a weighted average,
        $(1-\alpha)I_1 + \alpha I_2$. Used for the translucent HUD card with $\alpha = 0.6$.
  \item[Anchor] A predefined box in a single-shot detector. BlazePalm refines anchors into palm
        boxes.
  \item[BGR] The channel order OpenCV uses: blue, green, red. MediaPipe expects RGB, hence the one
        explicit conversion per frame.
  \item[BlazePalm] The stage-1 detector. Finds palm boxes, not fingers --- palms are rigid, so they
        are easier to detect robustly.
  \item[Cross product] A vector operation producing a vector perpendicular to two inputs. Its
        $z$-component here gives the signed area of the palm triangle, whose sign reveals which side
        of the hand faces the camera.
  \item[Debouncing] Requiring several consecutive agreeing observations before accepting a state
        change. Used on gestures: 4 of the last 5 frames.
  \item[Duck typing] Accepting any object that has the attributes you use, regardless of its class.
        This is what makes \code{parse\_hand\_data} testable with a fake result object.
  \item[EMA] Exponential moving average: a weighted blend of the current measurement and the previous
        smoothed value, with geometrically decaying memory. Used on the index fingertip with
        $\alpha = 0.65$.
  \item[Euclidean distance] Straight-line distance,
        $d = \sqrt{\Delta x^2 + \Delta y^2}$. The basis of the thumb rule.
  \item[Frame] One still image from a video stream.
  \item[FPS] Frames per second --- the loop's throughput. Averaged over a 30-frame window.
  \item[HUD] Head-up display: the text overlay showing FPS, finger count, gesture and orientation.
  \item[Inference] Running a trained model on new input. Training is not part of this project.
  \item[Landmark] One of the 21 named 3D points the model predicts per hand.
  \item[MCP / PIP / DIP / TIP] Knuckle, first finger joint, second finger joint, fingertip.
  \item[Normalized coordinates] Positions in $[0,1]$ relative to the image size, so the model is
        resolution-independent.
  \item[Palm normal] The vector perpendicular to the palm plane; its $z$-component's sign says
        whether the palm or the back faces the camera.
  \item[ROI] Region of interest --- the rectangular slice of an image being processed, here the HUD
        card area.
  \item[Single-pass caching] Computing derived telemetry once per frame and reusing it, rather than
        recomputing per call.
  \item[Static image mode] MediaPipe configuration that treats each frame as an unrelated still.
        Required for \code{--image} mode.
  \item[Tensor] A multi-dimensional array of numbers. A frame is a 3D tensor of shape
        $(H, W, 3)$.
\end{description}

\section{Constants and defaults}

\begin{center}
\rowcolors{2}{slatebg}{white}
\begin{tabularx}{\textwidth}{@{}l l l X@{}}
\toprule
\rowcolor{white}
\textbf{Name} & \textbf{Value} & \textbf{Where} & \textbf{Meaning} \\
\midrule
\code{min\_detection\_confidence} & 0.7 & tracker & Confidence to accept a \emph{new} hand. \\
\code{min\_tracking\_confidence} & 0.6 & tracker & Confidence to keep \emph{following} a hand. \\
\code{max\_num\_hands} & 2 & tracker & Hands searched for simultaneously. \\
\code{EXPECTED\_LANDMARKS} & 21 & logic & Validity gate for a landmark list. \\
\code{FINGER\_PAIRS} & \code{(8,6),(12,10),(16,14),(20,18)} & logic & Tip and PIP per long finger. \\
\code{THUMB\_TIP} / \code{THUMB\_IP} / \code{PINKY\_MCP} & 4 / 3 / 17 & logic & The thumb rule's three
  points. \\
\code{INDEX\_TIP} & 8 & logic & The landmark carrying the smoothed marker. \\
\code{alpha} (EMA) & 0.65 & logic & Smoothing weight on the newest sample. \\
\code{window} / \code{min\_agree} & 5 / 4 & logic & Debouncer window and required agreement. \\
\code{HUD\_ALPHA} & 0.6 & tracker & Overlay weight in the HUD blend. \\
\code{FPS\_WINDOW} & 30 & tracker & Frames averaged for the FPS readout. \\
DEFAULT\_WIDTH / HEIGHT & 640 / 480 & tracker & Live capture size. \\
\bottomrule
\end{tabularx}
\end{center}

\section{Dependency reference}

\begin{center}
\rowcolors{2}{slatebg}{white}
\begin{tabularx}{\textwidth}{@{}l l X@{}}
\toprule
\rowcolor{white}
\textbf{Package} & \textbf{Pin} & \textbf{Why this exact version} \\
\midrule
\code{mediapipe} & \code{0.10.14} & Last line with the legacy \code{mp.solutions.hands} API; removed
  in 0.10.30+. Load-bearing --- do not upgrade. \\
\code{opencv-contrib-python} & \code{4.11.0.86} & The superset build MediaPipe depends on; avoids two
  competing \code{cv2} installs. \\
\code{numpy} & \code{<2} & MediaPipe 0.10.x wheels are compiled against the NumPy 1.x ABI. \\
Python & 3.9--3.12 & 3.13 has no wheel for MediaPipe 0.10.14. Tested on 3.12. \\
\bottomrule
\end{tabularx}
\end{center}

\section{Files in the project}

\begin{center}
\rowcolors{2}{slatebg}{white}
\begin{tabularx}{\textwidth}{@{}l X@{}}
\toprule
\rowcolor{white}
\textbf{File} & \textbf{Purpose} \\
\midrule
\code{hand\_logic.py} & Pure logic: parsing, finger and thumb rules, palm normal, EMA, gestures,
  debouncing. Standard library only. \\
\code{hand\_tracker.py} & Application: \code{HandTracker}, capture loop, HUD, CLI, camera fallback,
  \code{--image} mode. \\
\code{tests/test\_hand\_logic.py} & Offline geometry and parse tests. No camera, no MediaPipe. \\
\code{tests/test\_hand\_tracker.py} & Cache and HUD tests; skips itself when cv2/mediapipe are
  absent. \\
\code{conftest.py} & Puts the project root on \code{sys.path} so \code{pytest} works as a bare
  console script. \\
\code{requirements.txt} & The three pins. \\
\code{run.bat} & Double-click launcher: uses the venv interpreter, passes arguments through. \\
\code{doc/build.bat} & Rebuilds this handbook (two \code{pdflatex} passes). \\
\code{.github/workflows/ci.yml} & Logic tests on 3.11/3.12, plus a full-dependency pipeline smoke
  test. \\
\bottomrule
\end{tabularx}
\end{center}

\vfill
\begin{center}
{\sffamily\color{slatetext}\small End of handbook.}
\end{center}


\end{document}
